\documentclass[11pt,letterpaper]{article}
\usepackage[T1]{fontenc}
\usepackage[utf8]{inputenc}
\usepackage{lmodern}
\usepackage{amsmath,amssymb,amsthm}
\usepackage[margin=1in]{geometry}
\usepackage{microtype}
\usepackage{booktabs,longtable,array}
\usepackage{xcolor}
\usepackage{listings}
\usepackage{enumitem}
\usepackage[hidelinks]{hyperref}
\usepackage{fancyhdr}
\usepackage{lastpage}
\setlength{\headheight}{14pt}
\pagestyle{fancy}
\fancyhf{}
\fancyhead[L]{\small Uniform successor evaluation for iPRA}
\fancyhead[R]{\small Research report}
\fancyfoot[C]{\small \thepage}
\setlength{\parskip}{4pt}
\setlength{\parindent}{0pt}
\setlist{itemsep=2pt,topsep=5pt}
\emergencystretch=2em
\newtheorem{theorem}{Theorem}[section]
\newtheorem{lemma}[theorem]{Lemma}
\newtheorem{proposition}[theorem]{Proposition}
\newtheorem{corollary}[theorem]{Corollary}
\theoremstyle{definition}
\newtheorem{definition}[theorem]{Definition}
\newtheorem{remark}[theorem]{Remark}
\lstset{basicstyle=\ttfamily\scriptsize,columns=fullflexible,
  keepspaces=true,breaklines=true,showstringspaces=false,
  frame=single,rulecolor=\color{black!20},xleftmargin=4pt,xrightmargin=4pt}
\hypersetup{pdftitle={One successor step for equality and computation},
  pdfauthor={Research prepared for Vladimir Reshetnikov},
  pdfsubject={Five-rule clausal presentation of intuitionistic primitive recursive arithmetic}}
\newcommand{\StepBytes}{2950}
\newcommand{\StepTokens}{829}
\newcommand{\StepAST}{902}
\newcommand{\ControlBytes}{3048}
\newcommand{\ControlTokens}{860}
\newcommand{\ControlAST}{949}
\newcommand{\StepSourceBytes}{3507}
\newcommand{\StepSourceGzip}{1274}
\begin{document}
\begin{titlepage}
\vspace*{1.0cm}
{\large Mathematical foundations \hfill 11 November 2026\par}
\vspace{1.2cm}
{\LARGE\bfseries One successor step for\\[4pt]
equality and computation\par}
\vspace{0.6cm}
{\Large A five-rule clausal presentation of\\
intuitionistic primitive recursive arithmetic\par}
\vspace{0.9cm}
{\large Prepared for Vladimir Reshetnikov\par}
\vspace{1.0cm}
\begin{abstract}
We give a simpler arithmetic inference interface for the current
\texttt{ConciseIPRA} proposal. On its already available zero-normal
iterator syntax, reflexivity and one uniform successor-evaluation rule
derive both computation and arbitrary equational transport. The three
syntactic alternatives of the current computational-equality rule are
replaced by one. The key is an explicit two-step derivation of equality
composition using a one-step identity iterator.

Together with the existing clausal logical interface, the complete
presentation has five primitive rule schemes: Use, Refl, Step, DE, and
DI. Its formulas are only equations and one two-variable binder.
We prove exact consequence equivalence with the native kernel and
therefore whole-domain bi-interpretability and equiconsistency with
full first-order iPRA. Both numerical composites have the identity
graph \(x=y\).

The fully validated standalone Python checker measures
\StepBytes\ normalized bytes, \StepTokens\ significant tokens, and
\StepAST\ AST nodes. A matched zero-normal computation/transport
control measures \ControlBytes, \ControlTokens, and \ControlAST,
respectively. The comparison counts normalization, syntax validation,
scope handling, clause storage, and the command-line interface.
The contribution is a local mathematical and implementation reduction;
no absolute minimum is asserted. The new equivalence is a written
proof supported by executable proof translations and boundary tests,
not a new Lean or Rocq verification.
\end{abstract}
\vfill
\textbf{Frozen comparison.}\quad
\texttt{VladimirReshetnikov/ProveIt}, commit\\
\texttt{6ca328039cafa1537d81ee6b67ee2dc54ffa5a4e}.\\
\textbf{Scope.}\quad Full intuitionistic first-order arithmetic,
including quantified formulas and arbitrary retained assumptions.
\end{titlepage}
\tableofcontents
\clearpage

\section{Result and comparison with the current proposal}

The useful simplification is the following single arithmetic inference.
Write \(\overline I\) for the canonical iterator constructor, which
returns its initializer at the literal counter zero and otherwise retains
an iterator node:
\[
 \boxed{\quad
 \frac{\overline I(c,a;u.h)=t\qquad I(Sc,a;u.h)=r}
      {r=N(h[t/u])}
 \quad\text{\rm Step}. \quad}
\]
Here \(N\) contracts zero-counter iterators recursively, and
substitution is capture-free. The rule reads its counter, initializer,
and body from the second premise. Its first premise is required to
have exactly the corresponding predecessor source. No arbitrary
equality template or occurrence path is supplied in the certificate.

In the frozen repository, \texttt{NativeEval.lean} implements the
six-rule computational-equality presentation with separate
\texttt{evalId}, \texttt{evalZero}, and \texttt{evalSucc} constructors
for its three Eval alternatives. The repository also contains
zero-normalization and initializer-only equality proposals.
The incoming \texttt{ipra\_clausal\_canonical.zip} already has a
five-rule clausal calculus, and
\texttt{ipra\_scalar\_uniqueness.zip} already replaces induction by
uniqueness of iteration. Both archives were inspected before selecting
this construction. The present proposal does not claim those earlier
ideas as new.

The new observation is that on zero-normal initialized iterators,
\emph{the successor alternative alone recovers the missing identity
alternative}. For \(J_s=I(S0,s;v.v)\), two Steps give
\[
 \frac{s=t\qquad
       \displaystyle\frac{s=r\qquad J_s=J_s}{J_s=r}\ {\rm Step}}
      {r=t}\ {\rm Step}.
 \tag{1}\label{eq:headline-pivot}
\]
This equality-composition rule works under arbitrary retained assumptions.
It does not depend on induction, numerical cases, or the logical binder.
Together with one-step application of a numerical template, it restores
the entire equality interface.

\subsection{Precise claims}

\begin{theorem}[Main result]\label{thm:main}
Let \(C_{\rm step}\) be the five-rule calculus specified in
Section~\ref{sec:five-pattern-clauses}. It is a faithful presentation of
a theory bi-interpretable with full first-order iPRA.
The interpretation domains are the whole numerical domains, interpreted
equality is actual equality, and both composite interpretations have
isomorphism graph \(x=y\). It is consequently equiconsistent with iPRA.

The new arithmetic inference interface is interderivable, under arbitrary
side assumptions, with normalized successor computation and
one-template equational transport. Replacing the latter pair by Refl
and Step therefore preserves the consequence relation, not merely the
set of true standard numerical equations.
\end{theorem}

The proof has three layers. First, the equality argument
(\ref{eq:headline-pivot}) recovers computation and transport.
Second, the clause interface recovers the native conditional-closure
binder and both case and progress introductions. A nonliteral canonical
term provably equal to zero handles the one delicate introduction case.
Third, the existing internal arithmetic bridge transfers the result to
iPRA. The first two layers are proved in full below; the exact inherited
bridge proofs are reproduced in the appendices with attribution.

\subsection{What the simplicity claim establishes}

There are several distinct comparisons, and they should not be
collapsed into one number:
\begin{itemize}
\item Against the current three-alternative Eval rule, Step has one
  arithmetic pattern. The zero equation is accounted for by the
  existing canonical syntax.
\item Against a matched zero-normal clausal calculus, the number of
  inference schemes remains five. Its explicit computation axiom and
  template-taking transport rule are replaced by Refl and Step.
  The certificate no longer carries either a computation triple or a
  transport template for those steps.
\item Against the incoming five-rule clausal canonical checker, the new
  complete checker is smaller in all three reported source coordinates.
  This comparison also changes the numerical representation, so the
  matched control is the cleaner measurement of the arithmetic change.
\end{itemize}

This report does not beat every historical source-size profile.
The repository contains an aggressively encoded 516-byte guarded
profile, and the incoming uniqueness package has an 866-byte fully
validated implementation of a different certificate protocol. The new
readable checker is not smaller than those programs. There is no
representation-independent theorem here about the globally simplest
axiomatization. The established result is the exact rule simplification
and its controlled complete-checker reduction.

\subsection{Attribution and proof status}

The scalar iterator, the \(D\) binder, clause contexts, the
witness-directed introduction, and zero-normalization are inherited
from the supplied repository. The contribution of this package is the
uniform Step reduction, its combination with the existing logical
interface, the written equivalence proof, and its executable
implementation and translations.

The baseline's Lean modules establish the inherited native/scalar/iPRA
bridge. We inspected those sources but did not compile a new Lean
formalization in this package. Passing Python certificates supplies
implementation evidence; the universal metatheorem is the written
proof, and the two kinds of evidence are kept separate.

% Includable mathematical note. Uses only standard amsmath/amssymb commands.
% Written for the pinned ProveIt baseline 6ca328039cafa1537d81ee6b67ee2dc54ffa5a4e.

\section{Uniform successor evaluation on zero-normal terms}
\label{sec:uniform-step-theory}

\subsection{Canonical syntax and what it costs}

Start with raw scalar terms
\[
 t ::= x\mid 0\mid S(t)\mid I(n,a;u.h),
\]
where $u$ binds only its occurrences in $h$. Define $N$ by structural
recursion, normalizing all immediate subterms and replacing
$I(0,a;u.h)$ by $a$ whenever the normalized counter is literally zero.
Only this zero rule is contracted. In particular,
$I(S0,a;u.h)$ is retained and is not definitionally $h[a/u]$.

The numerical terms of the new calculus are the fixed points of $N$.
Every displayed substitution means ordinary capture-free substitution
followed by $N$; write $h\langle t\rangle=N(h[t/u])$. Normalization applies
also under all term and formula binders. On formulas it changes numerical
terms only. Freshness is computed in the usual free-variable sense.

The elementary coherence equations are
\begin{align*}
 N(N(t))&=N(t),\\
 N(t[\sigma])&=N(N(t)[N\circ\sigma]),\\
 (t\langle\sigma\rangle)\langle\tau\rangle
 &=t\langle x\mapsto\sigma(x)\langle\tau\rangle\rangle.
\end{align*}
They follow by structural induction, with the binder cases using lifted
substitutions. Normalization commutes with renaming and introduces no
free variables. The contraction procedure and normalized substitution
belong to the trusted specification; deleting the zero computation from
the inference rules does not delete this implementation cost.

\subsection{One evaluation pattern suffices for equality and computation}

Use reflexivity and just the following evaluation inference:
\[
 \frac{\Gamma\vdash I(c,a;u.h)=t
       \qquad
       \Gamma\vdash I(Sc,a;u.h)=r}
      {\Gamma\vdash r=h\langle t\rangle}.
 \tag{Step}\label{rule:uniform-step}
\]
All terms in the display are interpreted by the canonical convention.
In particular, at $c=0$ the first source is simply $a$. The same
$c,a,h$ must occur in both sources. There is one case of this rule, no
comparison of a term with an arbitrary context and no search for an
occurrence to replace.

\paragraph{Successor computation.}
Take reflexivity as both premises of Step. Its conclusion is precisely
\[
 I(Sc,a;u.h)=h\langle I(c,a;u.h)\rangle.
 \tag{Succ}
\]
Zero computation is already the definition of the canonical constructor.

\paragraph{Equality pivot.}
For an arbitrary canonical term $s$, abbreviate
\[
 J_s=I(S0,s;v.v).
\]
Given $p:\Gamma\vdash s=t$ and $q:\Gamma\vdash s=r$, Step at
$c=0,a=s,h(v)=v$, with major premise $q$ and minor premise the
reflexivity of $J_s$, proves $J_s=r$. A second instance of the same
Step, now with premises $p$ and $J_s=r$, proves $r=t$. Thus
\[
 \frac{\Gamma\vdash s=t\qquad\Gamma\vdash s=r}
      {\Gamma\vdash r=t}
 \tag{Pivot}\label{rule:pivot}
\]
is a derived rule using two Steps and one auxiliary reflexivity.
This derivation is valid at every $s$, including literal zero.

\paragraph{Symmetry and transitivity.}
Pivot applied to $s=s$ and $s=t$ gives $t=s$. Given $s=t$ and $t=r$,
first obtain $t=s$ by symmetry and then apply Pivot with common source
$t$ to obtain $s=r$. Hence all equivalence-relation laws of equality
are available.

\paragraph{Congruence in every fixed term template.}
Fix a canonical body $h(v)$, allowing passive parameters and nested
iterators, and put $K=I(S0,s;v.h(v))$. Step with the reflexivities of
$s$ and $K$ proves $K=h\langle s\rangle$. Step with $s=t$ and the
reflexivity of $K$ proves $K=h\langle t\rangle$. Pivot gives
\[
 \Gamma\vdash h\langle s\rangle=h\langle t\rangle.
\]
Combining this with $h\langle s\rangle=r$ by Pivot gives the
one-template crossed transport rule
\[
 \frac{\Gamma\vdash s=t\qquad
       \Gamma\vdash h\langle s\rangle=r}
      {\Gamma\vdash r=h\langle t\rangle}.
 \tag{Tr}\label{rule:recovered-tr}
\]
Two instances give transport in an equation template. None of these
arguments uses equation induction, numerical case distinction,
successor injectivity, a logical connective, or a structural term
substitution theorem for whole derivations.

\paragraph{Reverse derivation.}
In the zero-normal native kernel, successor computation gives
$I(Sc,a;h)=h\langle I(c,a;h)\rangle$. The two premises of Step,
symmetry, transitivity, and term congruence therefore imply
$r=h\langle t\rangle$. Consequently Refl and Step are interderivable,
with arbitrary side assumptions, with the normalized successor
computation and the old one-template equality rule. This is an exact
consequence-equivalence theorem, not merely a standard-model soundness
observation.

\subsection{A normal zero witness for the single introduction}

The familiar single-introduction convention takes progress at a literal
equation with literal witness $0$, and cases otherwise. Its old alternate
zero witness $I(0,0;u.0)$ is no longer available: it is definitionally
zero in the canonical grammar. Instead use
\[
 q:=I(S0,0;u.u).
\]
This is a closed canonical term which is not the literal $0$.
Successor computation proves $q=0$; equivalently one Step on the
reflexivities of $0$ and $q$ proves it directly. This syntactic distinction
is the only property of the alternate witness needed below.

\section{A five-pattern clausal calculus}
\label{sec:five-pattern-clauses}

This subsection gives the full logical interface. It is independent of
whether a checker physically stores named variables, de Bruijn indices,
or an equivalent scoped encoding.

\subsection{Formulas and retained clauses}

Formulas have only the constructors
\[
 A ::= (s=t)\mid D_{x,y}B(x,y).
\]
The variables $x,y$ are bound in $B$. The ordinary interpretation is
\[
 \bigl(D_{x,y}B\bigr)^*
 :=\exists x\bigl(B(x,0)^*\to\forall y\,B(x,y)^*\bigr).
\]
The designated false formula is $0=S0$.
A retained clause is written
\[
 [A\,;\,y.B(y)]
\]
and has ordinary reading $A^*\to\forall y\,B(y)^*$; its variable $y$
is bound in $B$ and does not occur free in $A$. A context is a finite
list of such clauses. Empty contexts are allowed. For a formula $A$,
its ordinary-assumption encoding is
\[
 \operatorname{Assm}(A):=[0=0\,;\,y.A],\qquad y\notin\operatorname{FV}(A).
\]
This is an abbreviation for a clause, not a sixth inference rule.

\subsection{Complete rule inventory}

The following five patterns, and no further logical or arithmetic rules,
are primitive:
\begin{align*}
 \text{Use:}\quad&
 \frac{[A;y.B]\in\Gamma\qquad\Gamma\vdash A}
      {\Gamma\vdash B\langle t\rangle};\\[2mm]
 \text{Refl:}\quad&\Gamma\vdash t=t;\\[2mm]
 \text{Step:}\quad&
 \frac{\Gamma\vdash I(c,a;u.h)=t\qquad
       \Gamma\vdash I(Sc,a;u.h)=r}
      {\Gamma\vdash r=h\langle t\rangle};\\[2mm]
 \text{DE:}\quad&
 \frac{\Gamma\vdash D_{x,y}B\qquad
       \Gamma,[B(b,0);y.B(b,y)]\vdash A}
      {\Gamma\vdash A};\\[2mm]
 \text{DI:}\quad&
 \frac{\Gamma,\operatorname{Assm}(B(t,\eta))
          \vdash B(t,Sz)}
      {\Gamma\vdash D_{x,y}B}.
\end{align*}
For DE, $b$ is fresh for the outer context, the conclusion, and the
passive parameters. For DI, $z$ is fresh for the outer context, the
witness $t$, and all passive parameters; moreover
\[
 \eta=\begin{cases}
 z,& B\text{ is a literal equation and }t\text{ is literally }0,\\
 0,&\text{otherwise}.
 \end{cases}
\]
Every term and formula substitution in these rules is normalized.
Fresh variables used in a premise enlarge the scope of the entire outer
context as well as the formulas. Omitting that weakening is unsound.

All five patterns have one conclusion. Refl is an axiom schema.
Step supplies both computation and equational inference. Use accesses
retained clauses and thus includes the ordinary hypothesis rule:
from $\operatorname{Assm}(A)\in\Gamma$, apply Use to Refl$(0)$.
The only induction condition appears in the syntactic choice of
$\eta$ at DI; no induction for compound matrices is permitted.

\subsection{Structural bookkeeping}

Context weakening and renaming of free variables are admissible by
induction on derivations. Renaming preserves canonicality, literal zero,
and the root constructor of a formula, so the DI mode is preserved.
Ordinary formula cut is also admissible: if
$\Gamma\vdash A$ and $\Gamma,\operatorname{Assm}(A)\vdash C$, then
$\Gamma\vdash C$. Replace each Use of that assumption clause by the
first derivation, weakened to the current eigenvariable scope; induction
on the second derivation completes the argument. Its antecedent proof
of $0=0$ can be discarded. Bound eigenvariables are renamed before the
replacement when necessary.

These arguments do not require a direct proof of arbitrary term
substitution for the syntax-directed DI. Such a substitution may change
a nonliteral witness into literal zero. General term substitution becomes
admissible after equivalence with the ordinary native kernel has been
proved; treating its preservation of the mode as automatic would be an
error.

\subsection{Conditional closures as formulas}

Define the closure abbreviation
\[
 \widehat Q_y B(y):=D_{x,y}B(y),\qquad x\notin\operatorname{FV}(B).
\]
Its ordinary reading is equivalent to $B(0)\to\forall yB(y)$ because
the numerical domain has the closed inhabitant $0$.

\paragraph{Closure elimination.}
From $\Gamma\vdash\widehat Q_yB$ and $\Gamma\vdash B(0)$, apply DE.
Its retained clause is $[B(0);y.B(y)]$ and is independent of the fresh
witness. Use at any target $t$, with the weakened base proof, gives
$B(t)$. Thus QE is derived.

\paragraph{Closing a retained canonical clause.}
From a retained clause $[B(0);y.B(y)]$, derive $\widehat Q_yB$ by DI
with dummy witness $q$. This is the case mode, since $q$ is not literal
zero. Under the local assumption $B(0)$, Use proves $B(Sz)$.
The same argument converts any derivation
$\Gamma,\operatorname{Assm}(B(0))\vdash B(Sz)$, with $z$ fresh,
into a derivation of $\widehat Q_yB$: take the displayed derivation
as the premise of DI at dummy witness $q$.

\paragraph{Equation progress.}
Given
\[
 \Gamma,\operatorname{Assm}(E(z))\vdash E(Sz),
 \qquad E\text{ a literal equation},
\]
apply DI to the dummy-witness matrix $E(y)$ at witness $0$.
It produces $\widehat Q_yE$. Thus the native progress rule is derived
with no equality reasoning or additional induction schema.

\subsection{Recovering every case introduction}

For a compound matrix, or for a witness other than literal zero, DI is
already the desired case rule. Only a literal equation $E(x,y)$ with
witness $0$ requires an argument. Suppose
\[
 \Gamma,\operatorname{Assm}(E(0,0))\vdash E(0,Sz).
\]
By closing the displayed derivation at the dummy witness $q$, obtain
$\widehat Q_y E(0,y)$. Apply DE to retain
$[E(0,0);y.E(0,y)]$ and introduce $D_{x,y}E(x,y)$ by DI at witness
$q$, which uses the case mode. In its premise assume $E(q,0)$.
Transport along $q=0$ gives $E(0,0)$; Use of the retained clause gives
$E(0,Sz)$; transport back along $0=q$ gives $E(q,Sz)$, as required.
The fresh eigenvariable from the outer DE does not occur in the final
conclusion. All transports are in literal equation templates and are
available from the arithmetic Pivot construction.

Hence the five-pattern calculus derives both original introduction
rules, including the otherwise missing atomic case at zero. This proof
also explains why normalizing the old nonliteral zero witness without
replacing it is insufficient.

\section{Exact equivalence and the inherited iPRA bridge}
\label{sec:uniform-step-equivalence}

Let $N_0$ denote the original native calculus on raw scalar terms, and
let $N_z$ be its zero-normal presentation, with ordinary Hyp, successor
computation, one-template transport, QE, DE, case DI, and atomic QI.
The pinned baseline proves
\[
 N_0:\Gamma\vdash A
 \quad\Longleftrightarrow\quad
 N_z:N\Gamma\vdash NA.
\]
Its formal source is \texttt{NativeZeroNormal.lean}; the normal-form
coherence and interpretation facts are in
\texttt{ZeroNormalization.lean}. The semantic fact that a raw term
is provably equal to its normal form under an iterator body uses
body congruence, which in that inherited development uses equation
induction. The new equality bootstrap above does not claim to avoid
that inherited requirement in the proof of the normalization theorem.

Define $T$ on the formulas of $N_z$ by preserving equality and $D$,
and replacing every $Q_yB$ by $\widehat Q_y(TB)$. Translate ordinary
hypotheses $A$ to $\operatorname{Assm}(TA)$. The translation $T$
commutes with normalized numerical substitution up to renaming bound
variables: prove this by structural induction on formulas, choosing
each dummy witness binder fresh for the substitution images.
The preceding derivations
translate every rule of $N_z$ into the five-pattern calculus $C$:
Hyp becomes Use with Refl$(0)$; successor computation and transport
are derived from Step; QE uses DE and Use; equation progress is DI
at a dummy zero witness; every case introduction is recovered as
above. In the translation of DE, the retained canonical clause derives
the translated hypothetical closure, so admissible formula cut supplies
the translated branch derivation.

For the reverse direction, interpret retained clauses as
$A^*\to\forall yB(y)^*$ in the ordinary scalar theory $T_{\rm eq}$.
Use is implication elimination followed by universal elimination.
Refl is equality reflexivity. Step follows from successor computation
and equality. DE is ordinary existential elimination. In case-mode DI,
coverage $n=0\lor\exists z\,n=Sz$ combines the base assumption and
the successor premise to prove the conditional closure. In progress-mode
DI, contextual induction for the literal equation supplies that closure.
Existential introduction then supplies the stated witness in each mode.
Thus the entire calculus is sound in $T_{\rm eq}$ with precisely the
same domain and actual equality.

The pinned ordinary/native equivalence and the rule translations yield,
for every finite list $\Gamma$ of native formulas and every native $A$,
\[
 N_0:\Gamma\vdash A
 \quad\Longleftrightarrow\quad
 C:\operatorname{Assm}(T(N\Gamma))\vdash T(NA).
 \tag{Faithfulness}
\]
Here the right-hand context means pointwise application to its entries.
The reverse implication uses that the ordinary reading of
$\operatorname{Assm}(B)$ is equivalent to $B$, and that the ordinary
reading of the dummy-existential closure is equivalent to $Q$.
Those equivalences are intuitionistically valid over an inhabited domain.

\paragraph{Bi-interpretability, rather than only equiconsistency.}
The inherited bridge interprets raw scalar terms as their fixed primitive
recursive definitions in iPRA, and compiles each fixed primitive recursive
definition back to a scalar term. Its round trips are provably the identity
on numerical values. Compose those maps with $N$ and $T$. The normalizing
map fixes variables and zero, preserves numerical values provably, and
leaves the whole numerical domain unchanged. The closure replacement has
no numerical effect. Consequently both composites retain the identity
graph $x=y$; equality remains actual equality, and no quotient or proper
subdomain is introduced.

On formulas, inclusion followed by $T\circ N$ is the identity on the
canonical Eq/$D$ language. The other round trip is provably equivalent
to the input formula by zero-normalization equivalence and
$Q_yB\leftrightarrow D_{x,y}B$ for a dummy $x$. Transporting the inherited
ordinary/iPRA formula round-trip theorems through Faithfulness gives both
full first-order round trips for $C$. In particular, the designated
contradiction $0=S0$ is preserved, so the resulting theories are
equiconsistent as well as bi-interpretable.

\paragraph{Boundary of the claim.}
The new mathematical deletion is the replacement of the multi-case
arithmetic equality/computation mechanism by Refl and one uniform Step.
Zero-normalization and clausal logical contexts are inherited design
ideas and their implementation costs must be counted. A numerical size
comparison is a separate claim about complete programs, not a theorem
that no smaller presentation exists. Any proof-assistant status must be
reported according to the files actually compiled; the written
five-pattern equivalence is not automatically a machine-checked theorem.
\section{Complete verifier and certificate contract}
\label{sec:implementation}

The primary module \texttt{code/checker.py} is the entire trusted
acceptance program. It imports only Python's standard \texttt{json} and
\texttt{sys} modules. It reads one JSON object, validates the closed goal,
reconstructs the proof conclusion, and accepts exactly when they agree.
Acceptance is exit status zero and the line \texttt{ACCEPT}. Rejection
has nonzero exit status. Python's resource limits can reject sufficiently
large valid inputs; no resource-bound completeness claim is made.

\subsection{Concrete syntax}

The public object has exactly two fields, \texttt{goal} and
\texttt{proof}. All term and formula trees are explicitly checked:
\begin{center}
\begin{tabular}{lll}
\toprule
Mathematical object & JSON form & Binding\\
\midrule
Zero & \texttt{0} & None\\
Successor & \texttt{['s', a]} & None\\
Iterator & \texttt{['i', n, a, h]} & One variable in \texttt{h}\\
Equation & \texttt{['=', a, b]} & None\\
Paired formula & \texttt{['d', B]} & Two variables in \texttt{B}\\
\bottomrule
\end{tabular}
\end{center}
The display uses single quotes for readability; JSON files use double
quotes. A positive integer denotes a bound-variable index, with
index \(1\) the innermost variable. In a \(D\) matrix, index \(1\)
denotes \(y\) and index \(2\) denotes \(x\). Negative integers are
rigid free-variable levels used only inside local proof scopes.
At free scope \(k\) and binder depth \(d\), an integer term must satisfy
\(-k\leq a\leq d\). The number zero always denotes zero, not a variable.
Python booleans are explicitly excluded by exact integer-type tests.

Canonical term validation forbids an iterator with literal zero
counter, recursively everywhere. The shared opening operation rebuilds
every visited iterator using the canonical constructor, so exposing a
zero counter during substitution immediately contracts it. The test is
on the rebuilt counter, not on its original raw syntax.

The proof constructors are as follows:
\begin{center}
\begin{tabular}{ll}
\toprule
Rule & Certificate form\\
\midrule
Use & \texttt{['u', clause-index, witness, antecedent-proof]}\\
Refl & \texttt{['r', term]}\\
Step & \texttt{['x', major-proof, minor-proof]}\\
DE & \texttt{['e', goal, paired-proof, branch-proof]}\\
DI & \texttt{['i', paired-goal, witness, branch-proof]}\\
\bottomrule
\end{tabular}
\end{center}

The context is an internal list of antecedent/body pairs; it is not
an unchecked public environment. A DE branch opens its witness at
the next fresh negative level. A DI branch uses the next fresh level
as its induction/case parameter. The conclusion and chosen witness are
validated in the old scope before either branch opens, which prevents
an eigenvariable from escaping. Ordinary local assumptions are stored
as clauses with antecedent \(0=0\).

\subsection{Why the opening operation needs no shift}

The use of negative levels is an implementation convention, not an
extra axiom. Every substitution image is locally closed: all its
positive indices are bound within that image itself, and all its
external variables have negative labels. Inserting such an image
under an additional binder therefore cannot capture an external
variable. The opening traversal increases its depth under the
iterator body and by two under \(D\), replaces the designated bound
index, and decrements larger indices. The substituted image is not
recursively opened a second time.

This argument would fail for arbitrary open positive-index images.
The validator prevents that input. The proof-reconstruction invariant is
that every returned formula is canonical and well-scoped in the current
scope, and every retained clause has a well-scoped antecedent and a
body with exactly one additional permitted bound variable. Each
constructor preserves that invariant. In particular, Step's
premises are explicitly matched as equations with the required
successor-iterator source.

\subsection{Measurements}

All source measurements count the entire acceptance module.
For comparison with the supplied repository, module, class, and
function docstrings are removed from the AST, then the source is
normalized by CPython 3.12.14 \texttt{ast.unparse}, followed by one LF.
Significant tokens exclude comments, whitespace, indentation, newline,
encoding, and end markers. AST nodes are counted after that same
docstring removal. Gzip uses maximum level \(9\) and timestamp zero;
it is informational, not an optimization objective.
For the delivered primary Python file itself, including its documentation
and original formatting, the corresponding figures are
\StepSourceBytes\ source bytes and \StepSourceGzip\ gzip-9 bytes.
Both raw and normalized compressed files are included.

\begin{center}
\small
\begin{tabular}{lrrrr}
\toprule
Complete checker & Bytes & Tokens & AST & Gzip 9\\
\midrule
Matched zero-normal Comp/Tr & 3048 & 860 & 949 & 1018\\
Matched raw three-case Eval & 3094 & 853 & 969 & 1026\\
Incoming canonical clausal & 3305 & 946 & 1072 & 1073\\
New uniform Step & 2950 & 829 & 902 & 1005\\
\bottomrule
\end{tabular}
\end{center}

The matched canonical control shares the terms, normalization,
opening operation, formula grammar, public fields, validation, clause
protocol, and logical rules with the new checker. Its arithmetic rules
are successor computation and one-template crossed transport.
Consequently that comparison isolates the arithmetic replacement.
The raw Eval control retains raw zero-counter iterators and the three
arithmetic Eval alternatives. Its comparison also includes the cost
of adding canonical syntax.

The secondary modules in \texttt{code/native/} keep the more general
explicit-sequent Eq/\(Q\)/\(D\) interface with a variable count and
public hypotheses. There the matched normalized control measures
3,118 bytes, 1,010 tokens, and 1,080 AST nodes; Refl/Step measures
3,029 bytes, 981 tokens, and 1,051 AST nodes. The pinned raw Eval
checker measures 3,039 bytes, 959 tokens, and 1,026 AST nodes.
Thus the change has a clear matched win, but its comparison with
raw Eval in that second interface is mixed. These modules are
supplementary evidence, not additional hidden dependencies of
the primary verifier.

\subsection{Verification evidence and its limits}

\begin{itemize}
\item Primary clausal suite: 376 standalone-process checks in ordinary and optimized Python; no failures.
\item Explicit-sequent suite: 38 positive certificates, 12 negative controls, and 38 complete proof-translation round trips; 100 command-line checks.
\item Independent opening audit: 2104 semantic comparisons, 13 targeted malformed/capture controls, and one positive smoke check.
\end{itemize}

The explicit-sequent suite supplies actual finite proof translations
in both directions between normalized computation/transport and
Refl/Step. The translators are untrusted producers: each output is
checked by the target's complete verifier. They cover every primitive
constructor of those calculi, and preserve the public variable count,
context, and goal.

The tests specifically address the difficult implementation boundaries:
an induction hypothesis that is actually used, a closure actually
used after \(D\)-elimination, nested term binders, normalization
exposed by substitution, arbitrary equation assumptions, malformed
indices, booleans, forged predecessor sources, and escaping
eigenvariables. An independently written numerical evaluator also
checks the opening operation on finite instances. Numerical tests
are supplementary and are not used to infer the nonstandard-model
interpretation theorem.

The script \texttt{reproduce.py} reruns the validation suites and
recomputes all source metrics and deterministic gzip files.
Its optional \texttt{--pdf} switch rebuilds the standalone
\texttt{report.tex} and compiles the PDF. The inherited reference
files are byte-authenticated against Git blob identifiers at the
frozen commit; \texttt{evidence/provenance.json} records that check.

\section{Full first-order interpretation and research boundary}
\label{sec:interpretation-summary}

The interpretation into ordinary logic uses the entire numerical
domain and interprets each native equality as actual equality.
The \(D\) binder and clauses have precisely the ordinary readings
specified above. Genuine ordinary falsity is interpreted as \(0=S0\).
The ordinary scalar realization explicitly includes the separation
axiom \((0=S0)\to\bot\); positive numerical equations alone do not
supply this ordinary falsity bridge.

For transparency, the complete inherited logical macro construction
is reproduced in Appendix~\ref{sec:choice-introduction}.
It derives implication, conjunction, disjunction, both quantifiers,
and explosion using only selectors, finite iteration, and the
native rules. It also proves the paired-formula round trip
\[
 D_{x,y}B
 \quad\dashv\vdash\quad
 \exists_*x\bigl(B(x,0)\Rightarrow_*\forall_*y\,B(x,y)\bigr).
\]
This is why the result concerns full intuitionistic first-order
arithmetic rather than just an equational or negative fragment.
No excluded middle for quantified formulas is introduced.

The arithmetic bridge is also internal. Its order matters:
first construct predecessor by scalar iteration and equation
induction; then prove faithful equality reflection; then obtain
quantifier-free induction; only then construct pairing and compile
general primitive recursion. The scalar update has one bound
state variable and no hidden recursion-counter argument.

For a target recursion
\[
 f(0,\vec p)=g(\vec p),\qquad
 f(Sn,\vec p)=h(n,f(n,\vec p),\vec p),
\]
the inherited pairing \(\pi\) and projections \(L,H\) give
\[
 V(n,\vec p)=I\bigl(n,\pi(0,g^*(\vec p));
     u.\pi(SL(u),h^*(L(u),H(u),\vec p))\bigr),
 \qquad f^*(n,\vec p)=H(V(n,\vec p)).
\]
Equation induction proves \(L(V(n,\vec p))=n\), and the projection
equations prove both defining equations for \(f^*\).
Each fixed finite primitive-recursive description is compiled
externally. There is no internally total universal evaluator over
arbitrary descriptions.

In the other direction, each fixed scalar iterator is an ordinary
primitive recursion ignoring its counter in the update.
Equation induction proves uniqueness for these fixed recursions,
so an external induction over descriptions proves both internal
numerical round trips. Normalization preserves numerical values,
and the clause translation does not change numerical terms.
This establishes the identity graph \(x=y\) on the whole domain
in both composites, including nonstandard elements.

\subsection{What remains open}

The five-rule inventory retains a syntax-directed case/progress
choice inside DI. It has not removed the substantive distinction
between numerical coverage and equation induction. Further merging
the logical rules would need to preserve that distinction and its
eigenvariable conditions.

The result also exchanges a larger equality interface for a modest
increase in some proof sizes: derived equality composition expands
to two Step nodes and an auxiliary reflexivity. Checker source
size, certificate length, and checking time are separate objectives.
The implementation measurements do not establish an optimum for
any of them.

The next useful formal-verification target is a deeply embedded
Step calculus and its equivalence with
\texttt{ConciseIPRA.NativeZeroNormal}. The present package provides
the necessary derivations and tested finite translations, but
does not label an uncompiled Lean sketch as a verified result.

\clearpage
\appendix
\section{Inherited proof material and attribution}

The next five appendices reproduce the exact mathematical source
fragments from \path{Research/certified_zero/reference/} at the
frozen commit. The comparison script verifies their Git blob hashes.
They are included under the repository's MIT-0 license, retained
as \path{reference/LICENSE-ProveIt.txt}. Their descriptions of
earlier contributions are relative to their original reports.
They do not assert priority for this package.

These fragments use the raw scalar notation and their own earlier
native rule presentation. Sections~\ref{sec:uniform-step-theory}
through~\ref{sec:uniform-step-equivalence} prove the equivalence that
allows their arithmetic bridge to be used here. In particular,
the fragments' full formula-transport rule is a recovered rule,
not an additional primitive of the five-rule calculus.


% New mathematical result: complementary introduction choices.
% Requires amsmath/amssymb/amsthm and theorem, lemma, proposition environments.
% The inherited arithmetic-to-iPRA bridge must be supplied separately.

\section{Complementary introductions with one local assumption}
\label{sec:choice-introduction}

The old paired introduction retains both the zero instance and an
induction instance when its matrix is an equation. Both together are
unnecessary. A certificate may choose either one, provided the induction
choice is restricted to a literal equation. This changes neither the
formula language nor the meaning of a formula. The omitted combination
is recovered by a short nested introduction--elimination construction.

\subsection{Complete native syntax and rule inventory}

The scalar numerical language is
\[
 t::=x\mid0\mid St\mid I(n,a;u.h).
\]
Only $h$ binds the state variable $u$; $n$ and $a$ do not.
The update may contain fixed numerical parameters. The computation
axioms in every retained context are
\[
 I(0,a;u.h)=a,\qquad
 I(Sn,a;u.h)=h[I(n,a;u.h)/u].
 \tag{I0, IS}\label{eq:choice-computations}
\]
The argument also applies unchanged to the full fixed-description PR
language with its projection, composition, and recursion equations.
No numerical variable ranges over function descriptions, and no uniform
evaluator is postulated.

Native formulas and reusable context clauses are
\[
 A::=(s=t)\mid\mathsf D_{x,y}A,\qquad
 \langle A;y.B(y)\rangle,\quad y\notin\operatorname{FV}(A).
\]
Bound variables are renamed before capture-free substitution.
The two variables $x,y$ in a paired formula are distinct.
An ordinary assumption $F$ is the constant clause
$\langle H;v.F\rangle$, where $v$ is absent from $F$ and $H$ is the
fixed computation theorem $I(0,0;u.0)=0$. Thus ordinary hypothesis
use is a clause Use preceded by this computation proof.

There are computation axioms and exactly the following inference
patterns:
\[
\frac{\Gamma\vdash A}{\Gamma\vdash B(t)}
\quad\mathsf{Use}
\quad(\langle A;y.B(y)\rangle\in\Gamma),
\]
\[
\frac{\Gamma\vdash\mathsf D_{x,y}B(x,y)\qquad
      \Gamma,\langle B(a,0);y.B(a,y)\rangle\vdash F}
     {\Gamma\vdash F}
\quad\mathsf{DE},
\qquad
\frac{\Gamma\vdash s=t\qquad\Gamma\vdash F(s)}
     {\Gamma\vdash F(t)}
\quad\mathsf{Tr},
\]
\[
\frac{\Gamma,B(t,\eta_\varepsilon(z))\vdash B(t,Sz)}
     {\Gamma\vdash\mathsf D_{x,y}B(x,y)}
\quad\mathsf{DI}_{\varepsilon},
\qquad
\eta_0(z)=0,\quad\eta_1(z)=z.
\tag{Choice DI}\label{eq:choice-di}
\]
Here $\varepsilon$ is a certificate choice in $\{0,1\}$, and
$\varepsilon=1$ is permitted only when $B$ has a literal equation
root. An equation beneath a binder is not a literal equation.
Call $\varepsilon=0$ the \emph{case choice} and $\varepsilon=1$
the \emph{progress choice}. Both retain exactly one ordinary
formula assumption. Compound matrices have only the case choice.

In DI the eigenparameter $z$ is fresh for $\Gamma$, the witness
$t$, and every fixed parameter of $B$. In DE the eigenparameter
$a$ is fresh for $\Gamma$, the conclusion $F$, and every fixed
parameter of $B$. In transport, $F(v)$ is any fixed native
formula template. Displayed local assumptions and clauses are
discharged. There are no primitive reflexivity, coverage,
separation, ordinary logical, or additional induction rules.

We call this calculus $\mathcal C$. The original comparison
calculus $\mathcal N$ has identical syntax, computation,
Use, DE, and transport, but its introduction is
\[
\frac{\Gamma,B(t,0),[B(t,z)]_{\mathrm{at}}\vdash B(t,Sz)}
     {\Gamma\vdash\mathsf D_{x,y}B(x,y)}.
\tag{Old DI}\label{eq:choice-old-di}
\]
The bracket supplies $B(t,z)$ exactly at a literal equation.

\subsection{Exact equivalence and the unused-witness construction}

\begin{lemma}[Structural admissibility]
\label{lem:choice-structural}
Weakening, exchange, contraction, renaming, ordinary-assumption cut,
and capture-free numerical substitution are admissible in $\mathcal C$.
\end{lemma}
\begin{proof}
Induct on derivations. Rename local eigenparameters before inserting
proofs or terms. Numerical substitution preserves a formula's root
constructor, so an allowed progress choice remains allowed.
The selected choice does not change under substitution or cut.
For ordinary-assumption cut, replace each use of its constant clause
by the supplied proof. Computation schemas are closed under
capture-free numerical substitution. These facts verify every case.
\end{proof}

\begin{theorem}[Identity-language equivalence]
\label{thm:choice-equivalence}
For every retained clause context $\Gamma$ and native formula $F$,
\[
 \Gamma\vdash_{\mathcal C}F
 \quad\Longleftrightarrow\quad
 \Gamma\vdash_{\mathcal N}F.
\]
The transformations are effective on finite derivations and leave
all numerical and formula syntax unchanged.
\end{theorem}
\begin{proof}
Each new introduction is an instance of the old one after weakening
by the omitted local assumption. For the case choice at an equation,
add the induction instance. For the progress choice, add the zero
instance. Compound case introduction is already the old rule.
Thus every new derivation translates to an old one.

For the reverse direction, transform an old derivation recursively.
Only its atomic DI case needs a construction. Suppose its matrix is
the equation $E(x,y)$ and its transformed premise is
\[
 \pi:\quad\Gamma,E(t,0),E(t,z)\vdash E(t,Sz).
\]
The outer new inference will be case DI with witness $t$. It is
therefore enough to prove $E(t,Sz)$ under
$\Gamma,E(t,0)$.

Rename the eigenvariable in $\pi$ to a fresh $v$. Choose a new
variable $u$ absent from $E(t,y)$ and all retained data. Under
$\Gamma,E(t,0)$, progress DI with witness zero introduces
\[
             \mathsf D_{u,y}E(t,y)
\]
from exactly the renamed premise $\pi[v/z]$.
The bound coordinate $u$ is deliberately unused.

Apply DE to this paired formula. Because $u$ is unused, the opened
clause is exactly
\[
               \langle E(t,0);y.E(t,y)\rangle;
\]
it contains no new existential witness. The retained ordinary
assumption $E(t,0)$ supplies its antecedent. Use at $Sz$ gives
$E(t,Sz)$, and DE closes. Finally the outer case DI introduces
the original $\mathsf D_{x,y}E(x,y)$ with witness $t$.
The original premise proof $\pi$ has been inserted once.

An old compound introduction already is the new case choice.
All other rule cases are unchanged. This completes the induction
on the original derivation.
\end{proof}

The unused inner coordinate is essential. Introducing
$\mathsf D_{x,y}E(x,y)$ internally and then eliminating it would
give an arbitrary fresh witness, not necessarily the original
$t$. The construction instead keeps $t$ as a fixed parameter
of $E(t,y)$ and binds only the unused coordinate $u$.

\begin{proposition}[Linear bound on inference nodes]
\label{prop:choice-node-bound}
Count each computation axiom, Use, DI, DE, and transport as one
inference node, expanding ordinary-assumption use as a computation
of $H$ followed by Use. If an old derivation has $N$ nodes and
$r$ atomic DI nodes, the displayed transformation has exactly
\[
                   N+5r\leq 6N
\]
nodes.
\end{proposition}
\begin{proof}
Each atomic old DI node is replaced by six nodes: outer case DI,
inner progress DI, DE, Use of the opened clause, Use of the
retained base assumption, and the computation of $H$ needed for
that last ordinary-assumption use. Its old premise subtree
occurs once. All other inference nodes are copied once.
\end{proof}

This counts inference-tree nodes. It is not a bound on serialized
JSON bytes, formula sizes, checker runtime, or the number of
consumed packet cells when hypothesis closure is implicit.
In the formula-context Q/D presentation, the corresponding
replacement has outer case DI, inner progress DI, DE, QE,
and two Hyp nodes, giving the same six-node count.

\subsection{Induction and numerical cases separate cleanly}

\begin{lemma}[Contextual equation induction]
\label{lem:choice-ei}
The progress choice, DE, and Use derive
\[
\frac{\Gamma\vdash E(0)\qquad
      \Gamma,E(z)\vdash E(Sz)}
     {\Gamma\vdash E(t)}
\tag{EI}\label{eq:choice-ei}
\]
for every equation $E$, with the usual freshness condition.
\end{lemma}
\begin{proof}
Choose an unused witness coordinate $u$. Progress DI introduces
$\mathsf D_{u,y}E(y)$ from the given step. DE opens the clause
$\langle E(0);y.E(y)\rangle$. Supply its antecedent using the given
base proof and use the clause at $t$. Close DE.
\end{proof}

\begin{lemma}[Cases, equality, and binder congruence]
\label{lem:choice-preinduction}
The case choice derives, for arbitrary native formulas,
\[
\frac{\Gamma\vdash F(0)\qquad\Gamma\vdash F(Sz)}
     {\Gamma\vdash F(t)}
\tag{NC}\label{eq:choice-nc}
\]
with the usual freshness condition. Equality reflexivity,
symmetry, transitivity, numerical congruence, and uniform
matrix interderivability under $\mathsf D$ are derivable
without using the progress choice.
\end{lemma}
\begin{proof}
For cases, use case DI with an unused witness coordinate and
matrix $F(y)$, ignoring its retained base in the successor
premise. DE opens $\langle F(0);y.F(y)\rangle$; the given zero
proof and Use at $t$ finish.

For reflexivity put $r=I(0,t;u.0)$. Use the computation equation
$r=t$ as both transport premises with motive $v=t$ to obtain
$t=t$. The usual explicit equality motives give symmetry,
transitivity, and numerical congruence.

For binder congruence, open the old paired formula at a fresh
witness $a$ and introduce the new one by case DI at $a$.
Convert its retained base to the old base, use the opened
old clause at the successor, and convert back. Ordinary cut
closes the construction, and the reverse conversion gives
the converse. No induction is used.
\end{proof}

The primitive progress choice is still restricted to equations.
Allowing it at arbitrary compound matrices would, by the
argument for \eqref{eq:choice-ei}, supply induction for
arbitrary formulas and change the intended theory.

\subsection{The entire ordinary logical interface}
\label{subsec:choice-logic}

Here is a predecessor-free reconstruction of the ordinary logical
operations. It makes explicit that the case choice supplies the
logical interface independently of progress induction.

Let $C(v,a,b):=I(v,a;u.b)$ with $u$ unused. Computation gives
$C(0,a,b)=a$ and $C(Sv,a,b)=b$. Define the gate on fixed finite
formula syntax by
\[
 G_v(s=t):=(C(v,0,s)=C(v,0,t)),\qquad
 G_v(\mathsf D_{x,y}B):=\mathsf D_{x,y}G_vB,
\]
renaming bound variables away from the control. Equality,
binder congruence, and external structural induction give
\[
 \vdash G_0A,\qquad G_{Sv}A\dashv\vdash A,\qquad
 A\dashv\vdash\mathsf D_{u,y}G_yA
 \quad(u,y\notin\operatorname{FV}(A)).
\tag{Gate and padding}\label{eq:choice-gate}
\]
For padding, case DI proves the forward direction. DE, the
zero-gate theorem, and Use at $1$ prove the reverse direction.

Define a selector between two arbitrary formulas by
\[
 M_v(s=t,u=w):=(C(v,s,u)=C(v,t,w))
\]
and, with aligned fresh paired variables, by
\[
\begin{aligned}
 M_v(\mathsf D_{x,y}A,\mathsf D_{x,y}B)
   &:=\mathsf D_{x,y}M_v(A,B),\\
 M_v(E,\mathsf D_{x,y}B)
   &:=\mathsf D_{x,y}M_v(G_yE,B),\\
 M_v(\mathsf D_{x,y}A,E)
   &:=\mathsf D_{x,y}M_v(A,G_yE),
\end{aligned}
\]
where $E$ is an equation. Every call decreases maximum input
binder depth. Padding and congruence prove
$M_0(A,B)\dashv\vdash A$ and
$M_{Sv}(A,B)\dashv\vdash B$.

Put $\delta(v)=C(v,1,0)$ and $p(z)=I(z,0;u.\delta(u))$.
Only the fixed computations $p(1)=1$ and $p(2)=0$ are needed.
The full logical macro inventory is
\[
\begin{aligned}
 A\Rightarrow_*B&:=\mathsf D_{u,z}M_z(A,B),\\
 A\wedge_*B&:=\mathsf D_{u,z}G_zM_{p(z)}(A,B),\\
 \exists_*x\,A(x)&:=\mathsf D_{x,y}G_yA(x),\\
 A\vee_*B&:=\exists_*x\,M_x(A,B),\\
 \forall_*x\,A(x)&:=
        A(0)\wedge_*\mathsf D_{u,z}G_zA(z),\\
 \bot_*&:=(0=1).
\end{aligned}
\tag{Logical macros}\label{eq:choice-macros}
\]
Newly bound variables are fresh; $u$ is unused.

Implication introduction obtains its assumption from the zero
selector and supplies the consequent at successors, then uses
case DI. Elimination opens its clause and uses it at $1$.
Existential introduction uses its witness and the successor
gate; elimination opens at a fresh witness, supplies the
zero-gate theorem, and uses the clause at $1$.

Disjunction introduction uses witnesses $0$ and $1$.
For elimination, discharge the two branch proofs as
$A\Rightarrow_*F$ and $B\Rightarrow_*F$. NC applied to
$M_z(A,B)\Rightarrow_*F$ gives this implication at every
control: the two endpoint cases follow from the branch
implications. Existential elimination provides $M_a(A,B)$
and the implication gives $F$.

Given $A$ and $B$, selector endpoints and NC give $M_v(A,B)$
at an arbitrary control. Substitute $v=p(Sz)$ and use the
successor gate for conjunction introduction. For projections,
open its zero-gated clause and specialize at $1$ and $2$;
the two fixed computations give $B$ and $A$, respectively.

These rules prove existential numerical coverage:
\[
     (x=0)\vee_*\exists_*v\,(x=Sv).
\tag{Coverage}\label{eq:choice-coverage}
\]
Use NC. Zero takes the left injection of reflexivity, and a
successor takes the right injection with the evident witness.
Coverage gives equality cases even when the retained context
contains $x$; its elimination adds $x=0$ or $x=Sv$ rather
than silently changing that retained context.

Universal introduction follows from a fresh-parameter proof
by substituting zero and successors and forming the two
conjuncts. For elimination, take its $A(0)$ conjunct and
open the paired conjunct. Its base is a zero-gate theorem.
Coverage of the target chooses the zero instance from
$A(0)$ or a successor instance from the opened clause
and gate removal.

Finally, under $0=1$, computation and transport give
$s=C(0,s,t)=C(1,s,t)=t$ for arbitrary numerical terms.
Every equation follows, and external induction on formula
syntax with case DI gives every native formula. This
supplies falsity elimination. No excluded middle for
quantified formulas has entered the construction.

\subsection{Whole-domain realizations and both formula composites}

For a standard first-order reading, let $T_{\rm eq}$ have
intuitionistic equality logic, the same numerical signature
and computation equations, numerical coverage, separation
$(0=1)\to\bot$, and contextual equation induction.
Its interpretation of the native syntax is
\[
 (s=t)^\circ=(s=t),\qquad
 (\mathsf D_{x,y}B)^\circ
   =\exists x\bigl(B^\circ(x,0)\to\forall y B^\circ(x,y)\bigr),
\]
\[
 \langle A;y.B\rangle^\circ
       =A^\circ\to\forall y B^\circ.
\]
The domain is the entire numerical domain and interpreted
equality is actual equality.

Case DI is sound in this realization by numerical coverage.
Progress DI is sound by contextual equation induction under
the conditional base. Use, DE, transport, and computation
have their displayed ordinary meanings. Conversely, the
native macros derive every ordinary logical rule,
\eqref{eq:choice-ei} derives equation induction, and
\eqref{eq:choice-coverage} derives coverage. Separation
translates to the identity implication
$(0=1)\Rightarrow_*(0=1)$.

For the essential paired formula composite,
\[
 \mathsf D_{x,y}B
 \dashv\vdash
 \exists_*x\bigl(B(x,0)\Rightarrow_*\forall_*y\,B(x,y)\bigr).
\tag{Paired composite}\label{eq:choice-paired-composite}
\]
Forwards, DE opens its clause at a fresh $a$; under $B(a,0)$
that clause supplies $B(a,y)$ at a fresh arbitrary $y$.
Introduce the universal, implication, and existential
macros and close DE. Backwards, eliminate the existential
at a fresh $a$. Introduce the paired formula with witness
$a$ by case DI: the retained base, supplied implication,
and universal elimination at $Sz$ give its successor premise.

The ordinary interpretations of the six logical macros
have their ordinary round trips. For conjunction, its true
base and the instances $1,2$ recover both operands;
the converse uses selector cases. The universal has all
successor instances in its paired conjunct and zero in
its first conjunct; coverage gives exactly the ordinary
universal. Implication, disjunction, and existential
quantification use their selector or gate endpoints.
Falsity uses the explicit ordinary separation bridge.
External induction on finite formulas, together with
\eqref{eq:choice-paired-composite}, gives both formula
composites. Clauses correspond in both directions to
$A\Rightarrow_*\forall_*yB$ by the derived logical rules.

The numerical composites are literal identities with
internally provable graph $x=y$. No standard-numeral
restriction, retract, quotient, or equality reflection
assumption is being inserted at this stage. The separately
supplied arithmetic bridge identifies $T_{\rm eq}$ with
full iPRA; the present change preserves that bridge
unchanged because Theorem~\ref{thm:choice-equivalence}
preserves every native sequent.

\subsection{What the theorem does and does not claim}

The new inference interface retains one ordinary assumption
at every paired introduction. It preserves the entire
consequence relation, not just equations or closed
theorems. Consequently the established full-iPRA
interpretations and equiconsistency transfer unchanged.

This result is different from retaining only progress
introduction and strengthening closure elimination.
Here the original elimination rule is retained, and
the case and progress introductions coexist. It is also
different from allowing arbitrary compound progress,
which would supply stronger induction.

A complete-checker comparison must still count the
certificate choice, its equation-root restriction,
substitution, context handling, computation, and the
public input/output interface. The mathematical theorem
alone is not a byte-count or minimality theorem.
The inference-node bound is distinct from source-size
measurements and from serialized certificate size.
The proof is a written universal proof; finite regression
tests do not constitute a Lean/Rocq formalization of it.

\clearpage
\section{Arithmetic strength and both ordinary interpretation composites}
\label{sec:ordinary-bridge}

We spell out the arithmetic bridge to avoid identifying standard numerical
agreement with bi-interpretability. First elaborate every fixed iterator
into an ordinary function symbol. At stage zero the symbols are $0,S$.
Given a fixed earlier-stage term $h(u,\vec p)$, introduce
$i_h(n,a,\vec p)$ with the two equations
\[
i_h(0,a,\vec p)=a,\qquad
i_h(Sn,a,\vec p)=h(i_h(n,a,\vec p),\vec p).
\]
Only finite descriptions occur; the domain contains no code for a
universal evaluator. Formal parameters are distinct when a description
is declared, while actual arguments can coincide.

Let $T_I$ be ordinary intuitionistic first-order logic with actual
equality and genuine falsity, these staged equations, contextual
induction for literal equations, and
\[
\forall x\,(x=0\vee\exists y\,x=Sy),
\qquad (0=S0)\to\bot.
\]
The second sentence is the falsity bridge. It must be stated: positive
numerical equations by themselves admit a singleton structure. Native
explosion from $0=S0$ and a translation of ordinary falsity to that
equation do not authorize silently dropping the bridge from $T_I$.

The native-to-ordinary translation is
\[
(s=t)^\circ=(s=t),\qquad
(\mathsf D_{x,y}B)^\circ=
\exists x\bigl(B^\circ(x,0)\to\forall y B^\circ(x,y)\bigr).
\]
A clause $\langle A;y.B\rangle$ becomes $A^\circ\to\forall yB^\circ$.
For the auxiliary closure used by the compact checkers,
$(\mathsf Q_yB)^\circ=B^\circ(0)\to\forall yB^\circ$.
The interpretation domain is all elements, defined by $x=x$;
equality is actual equality.

Case DI is sound because numerical coverage combines its retained base
with its successor proof. Progress DI is sound by literal-equation
induction under a temporarily assumed base. Use, DE and transport are
ordinary logical inferences under the translations. Conversely the
derived logical operations, equation induction, cases and explosion in
Section~\ref{sec:choice-introduction} recover the rules and axioms of
$T_I$. An auxiliary closure is interderivable with
$\mathsf D_{v,y}B(y)$ for unused $v$: in one direction use case DI with
the closure and its base; in the other eliminate the paired formula
and return its supplied closure. Thus formulas and context clauses have
the stated meanings in both presentations.

\subsection{The arithmetic dependency order}

The scalar update has no recursion-counter argument. We do not conceal
one in an encoding. The inherited internal proofs, reproduced in the
appendices, establish the required arithmetic in this order:
\begin{enumerate}
\item Using equation induction and numerical cases, construct a fixed
scalar predecessor term and prove $P0=0$ and $P(Sn)=n$.
Existential coverage then gives $x=0\vee x=S(Px)$: in the successor
case $x=Sy$, apply $P$ to obtain $Px=y$ and substitute.
\item Construct truncated subtraction and a numerical characteristic
of actual equality. Prove its reflection on the entire original domain.
\item Compile each fixed quantifier-free formula into its characteristic
equation and recover every contextual quantifier-free induction instance.
\item With that induction available, construct pairing and projections
by a bounded square-root search and compile every fixed general
primitive-recursive description into a scalar iterator term.
\end{enumerate}
No pairing or faithful equality detector is assumed to obtain
predecessor. No appeal to standard-model correctness replaces an
internal equality. Every induction fixes its passive parameters;
diagonal substitution occurs only into a previously proved identity.

For a target recursion
\[
f(0,\vec p)=g(\vec p),\qquad
f(Sn,\vec p)=h(n,f(n,\vec p),\vec p),
\]
let $\pi,L,H$ be the proved pairing and projections. The compiler uses
\[
V(n,\vec p)=I\bigl(n,\pi(0,g^*(\vec p));
u.\pi(SL(u),h^*(L(u),H(u),\vec p))\bigr),
\qquad f^*(n,\vec p)=H(V(n,\vec p)).
\]
Equation induction proves $L(V(n,\vec p))=n$; the two projection
identities then prove the target defining equations. The construction
is external recursion on each finite PR description, not an internal
quantification over arbitrary functions.

\begin{theorem}[Whole-domain bi-interpretability and equiconsistency]
The ordinary realization $T_I$ and full first-order iPRA are
bi-interpretable. In both theories $x=y$ defines the isomorphism from
the original domain to its double interpretation. The native Choice-DI
calculus is therefore a faithful proof presentation of a theory
bi-interpretable and equiconsistent with iPRA.
\end{theorem}
\begin{proof}
Interpret an iterator in iPRA by primitive recursion that ignores its
counter argument in the update. iPRA proves its equations, literal
equation induction, coverage and separation. In the reverse direction
use the scalar compiler. The predecessor and equality-reflection proofs
supply the successor axioms and every quantifier-free induction instance.
All ordinary logical rules are preserved.

For two iterators with the same seed and provably equal updates,
equation induction on the equality of their values proves uniqueness at
every counter, with all parameters fixed. The corresponding argument
applies to ordinary primitive recursion. Induct externally on a finite
description: its double translation satisfies its original defining
equations, so internal uniqueness identifies the two functions. This
also proves coherence when substitution changes the staged syntactic
description of an update.

The domain and equality are unchanged; hence the identity graph is
internally total, bijective and preserves every interpreted operation.
The displayed ordinary meanings of the logical macros and the paired
binder give both formula composites by external induction on finite
formula syntax. Clauses are similarly equivalent to their conditional
universal meanings. The explicit falsity bridge preserves contradiction
in both directions. Effective translations of finite proofs give
equiconsistency, without a claim that either theory proves its own
consistency.
\end{proof}

\clearpage
% Adapted with attribution from the pinned unparenthesized one-hole archive.
\section{Internal scalar predecessor bootstrap}
\label{sec:hybrid-bootstrap}

This scalar bootstrap is inherited from the supplied scalar development~\cite{erasure} and is reproduced to keep the scalar checker self-contained. It is not a new contribution of this paper. The proof uses a factor map fixing odd states and a saturated-iteration construction of predecessor.

The available arithmetic terms are $0$, $S$, and
\[
 I(n,a;u.h),\qquad
 I(0,a;u.h)=a,\qquad
 I(Sn,a;u.h)=h[I(n,a;u.h)/u].
\]
The update binds only $u$.  All other variables in it are fixed numerical
parameters.  We use equality reasoning, the already reconstructed
zero/successor cases, and induction on a single literal equation.
Every induction below keeps its other parameters fixed.  Predecessor,
successor injectivity, pairing, and induction on arbitrary
quantifier-free formulas are not assumptions of the argument.

\subsection{Elementary terms and two finite cycles}

Write $1=S0$ and use larger numerals as successor abbreviations.  Define
\[
 C(n,a,b)=I(n,a;u.b),\qquad
 Z(n)=C(n,1,0),\qquad
 a+b=I(b,a;u.Su).
\]
The update variable of $C$ is absent from its body.  Consequently
\[
 C(0,a,b)=a,\qquad C(Sn,a,b)=b.
\]
Numerical cases show that $Z(n)$ is $0$ or $1$.
The usual commutative-monoid equations for addition follow by literal
equation induction.  For example, $0+b=b$ and $Sa+b=S(a+b)$ are
proved by induction on $b$; these yield commutativity by a further
equation induction, and associativity follows by induction on the last
summand.  We use these equations to reorder finite sums.  The notation
$2a$ means $a+a$ and $3a$ means $a+a+a$; no multiplication operation
is needed for this bootstrap.

Put
\[
 N(n)=I(n,0;u.Z(u)),\qquad
 j(u)=Z(u)+2N(u),\qquad
 M(n)=I(n,0;u.j(u)).
\]
The range of $N$ is contained in $\{0,1\}$: its zero value is $0$,
and its successor value is a value of $Z$.  This uses numerical cases
on its counter, not induction on a disjunction.  Hence
\begin{equation}
 N(n)+Z(N(n))=1,\qquad
 Z(Z(N(n)))=N(n),\qquad N(n+2)=N(n).
 \label{eq:hybrid-parity}
\end{equation}
Here the last equation follows from two successor computations and
the preceding equation.  Literal-equation inductions give
\begin{equation}
 N(2n)=0,\qquad N(2n+1)=1,\qquad N(3n)=N(n).
 \label{eq:hybrid-parity-scaling}
\end{equation}
For the first equation the successor step adds two to the argument of
$N$; the second follows by one computation.  In the third, increasing
$n$ adds three to the left argument and therefore applies $Z$ once to
both sides of the induction equation.

Finite computations now give
\[
 j(0)=1,\qquad j(1)=2,\qquad j(2)=0.
\]
The proof below needs just these three values and
$M(Sn)=j(M(n))$; it does not assume a division algorithm or a theorem
characterizing $M$ as remainder modulo three.

\subsection{A ramp with a global scaling identity}

Define
\[
 G(n)=I(n,0;u.u+M(u)+1).
\]
The central invariant is
\begin{equation}
 M(G(n))=N(n).
 \label{eq:hybrid-ramp-invariant}
\end{equation}
Its zero case is immediate.  For its successor step, retain the
single induction equation and split the Boolean value $N(n)$.
If it is $0$, the next $G$ state is $G(n)+1$; its $M$ value is
$j(0)=1$, equal to $N(Sn)$.  If it is $1$, the next $G$ state
is $G(n)+2$; two computations of $M$ change its value from $1$
to $2$ and then to $0$, again equal to $N(Sn)$.  This proves
the invariant by literal-equation induction.

It follows that
\[
 G(Sn)=G(n)+N(n)+1,\qquad G(n+2)=G(n)+3.
\]
Equation induction therefore yields
\begin{equation}
 G(2n)=3n.
 \label{eq:hybrid-ramp-double}
\end{equation}
A separate literal-equation induction proves
\begin{equation}
 G(3n)=3G(n)+N(n).
 \label{eq:hybrid-ramp-triple}
\end{equation}
For completeness, its successor calculation is
\begin{align*}
 G(3Sn)
 &=G(3n)+N(3n)+4
   =3G(n)+2N(n)+4,\\
 3G(Sn)+N(Sn)
 &=3G(n)+3N(n)+3+Z(N(n)).
\end{align*}
The expressions agree in both Boolean cases for $N(n)$.
The base is $0=0$.  In particular,
\eqref{eq:hybrid-ramp-triple} was proved internally; it was not
inferred from an intended floor-function interpretation of $G$.

Set
\[
 H(u)=C(N(u),G(u),u).
\]
The crucial choice is to retain $u$ when its parity is $1$.
The preceding equations give
\begin{equation}
 H(2u)=3u,\qquad H(1)=1,\qquad H(3u)=3H(u).
 \label{eq:hybrid-factor-map}
\end{equation}
To check the last identity, use $N(3u)=N(u)$ and split the Boolean
$N(u)$.  At $0$, both sides are $3G(u)$ by
\eqref{eq:hybrid-ramp-triple}; at $1$, both are $3u$.
This argument does not require that an arbitrary number has already
been represented as $2v$ or $2v+1$.

\subsection{A comparator from scalar iteration}

For every fixed update $h$, equation induction on $n$, keeping $m,a$
and all parameters fixed, proves the semigroup identity
\[
 I(m+n,a;h)=I(n,I(m,a;h);h).
\]
The base is the zero computation and the step applies $h$ to the
induction equation. This identity will also be used in the later
pairing construction.


For a fixed update $h$, equation induction on $y$, with its initializer
$a$ and all parameters of $h$ fixed, proves the rotation identity
\begin{equation}
 I(Sy,a;h)=I(y,h(a);h).
 \label{eq:hybrid-rotation}
\end{equation}
Both sides have base $h(a)$, and their successor equations apply the
same $h$ to the previous equality.  Here and below $I(y,a;h)$ is
an abbreviation for one fixed finite scalar term, not quantification
over functions.

By \eqref{eq:hybrid-factor-map}, another equation induction gives
\begin{equation}
 I(y,3a;H)=3I(y,a;H).
 \label{eq:hybrid-iterator-scaling}
\end{equation}
Its base is $3a=3a$.  In the step the fixed-parameter induction
equation rewrites the new left side as
\[
 H(3I(y,a;H))=3H(I(y,a;H)).
\]
The initializer parameter stays $a$ throughout this induction.
Equation induction also gives $I(y,1;H)=1$, since $H(1)=1$.

Define
\[
 E(x)=I(x,1;u.2u),\qquad
 L(x,y)=N\bigl(I(y,E(x);H)\bigr).
\]
Then the following entire recursion specification is provable:
\begin{equation}
 L(0,y)=1,\qquad L(Sx,0)=0,\qquad L(Sx,Sy)=L(x,y).
 \label{eq:hybrid-comparator}
\end{equation}
The first equation uses $E(0)=1$ and the fixed-point identity for
$H$.  The second uses $E(Sx)=2E(x)$ and $N(2a)=0$.
For the third,
\begin{align*}
 L(Sx,Sy)
 &=N\bigl(I(Sy,2E(x);H)\bigr)\\
 &=N\bigl(I(y,H(2E(x));H)\bigr)\\
 &=N\bigl(I(y,3E(x);H)\bigr)\\
 &=N\bigl(3I(y,E(x);H)\bigr)\\
 &=N\bigl(I(y,E(x);H)\bigr)=L(x,y).
\end{align*}
The successive steps use rotation, $H(2a)=3a$, iterator scaling,
and $N(3a)=N(a)$.  Thus $L$ has the comparison equations at all
arguments internally.  Its standard-model meaning as the indicator
of $x\le y$ is motivation, not a premise of the following argument.

\subsection{Saturated iteration gives predecessor}

We now need only the three equations
\eqref{eq:hybrid-comparator}.  Define
\[
 B_m(u)=C(L(Su,m),u,Su),\qquad
 W(t,m)=I(t,0;u.B_m(u)).
\]
The subscript in $B_m$ records a fixed parameter; $B_m$ is an open
term abbreviation, not a function variable.

Numerical cases on $e$ prove the selector equation
\[
 C(e,Sa,Sb)=S C(e,a,b).
\]
Equivalently one may use equation induction on $e$, whose induction
hypothesis is unused.  Therefore
\begin{align}
 B_{Sm}(Su)
 &=C(L(SSu,Sm),Su,SSu)\notag\\
 &=C(L(Su,m),Su,SSu)=S B_m(u),
 \label{eq:hybrid-conjugacy}\\
 B_{Sm}(0)&=C(L(S0,Sm),0,S0)=C(L(0,m),0,S0)=S0.
 \label{eq:hybrid-conjugacy-base}
\end{align}

Literal-equation induction on $t$, with $m$ fixed, proves
\begin{equation}
 W(St,Sm)=S W(t,m).
 \label{eq:hybrid-saturated-shift}
\end{equation}
The base is \eqref{eq:hybrid-conjugacy-base}.  The successor step is
\begin{align*}
 W(SSt,Sm)
 &=B_{Sm}(W(St,Sm))\\
 &=B_{Sm}(S W(t,m))\\
 &=S B_m(W(t,m))=S W(St,m).
\end{align*}
Only the second line uses the current induction hypothesis, at the
same fixed $m$.

Two subsequent equation inductions give
\begin{equation}
 L(Sn,n)=0,\qquad W(n,n)=n.
 \label{eq:hybrid-diagonals}
\end{equation}
The first uses the comparator recurrence.  The second uses
\eqref{eq:hybrid-saturated-shift} in its successor step.  The
substitution $t=m=n$ is made into that already proved universally
parameterized equation; no parameter of an active induction hypothesis
is changed.  Consequently
\begin{equation}
 W(Sn,n)=B_n(n)=C(L(Sn,n),n,Sn)=n.
 \label{eq:hybrid-extra-step}
\end{equation}

Finally define the finite scalar term
\begin{equation}
 P(m)=I\bigl(m,0;u.C(L(SSu,m),u,Su)\bigr).
 \label{eq:hybrid-predecessor-term}
\end{equation}
Its zero equation is immediate.  To justify its successor equation
without assuming a separate extensionality rule for iterator bodies,
put
\[
 V(t,m)=I\bigl(t,0;u.C(L(SSu,Sm),u,Su)\bigr).
\]
The comparator recurrence identifies this update pointwise with
$B_m(u)$.  Equation induction on $t$, with $m$ fixed, gives
$V(t,m)=W(t,m)$: their base values are $0$, and the successor step
uses the pointwise update equation and then the induction equation.
Now \eqref{eq:hybrid-extra-step} yields
\[
 \boxed{\qquad P(0)=0,\qquad P(Sm)=V(Sm,m)=W(Sm,m)=m.\qquad}
\]

\subsection{Dependency boundary}

All terms in this section are built from $0,S,I$ by finite substitution.
All inductions have a single equation as their invariant.  Numerical
cases and equality transport are the already reconstructed logical
resources.  Successor injectivity may now be obtained by applying $P$
to an assumed equation $Sx=Sy$.  Equality tests, reflection of
quantifier-free formulas, pairing, and the recovery of general
primitive recursion can then be developed in the established order.
They were not used in order to manufacture predecessor.

This argument is a written internal derivation in the stated theory.
The finite implementation checks do not constitute a proof-assistant
formalization of this entire metatheorem.

\clearpage
% Standalone section.  Requires only amsmath and amssymb in the parent document.
% No theorem environments, macros, or bibliography keys are assumed.
% All arithmetic claims in this section are proved here.

\section{From literal-equation induction to full intuitionistic PRA}
\label{sec:eic-completion}\label{app:arithmetic}

This section supplies an unconditional completion theorem. It is adapted
with attribution from the current forward-left draft~\cite{forward}. In the
scalar application used here, every function occurring in the proof is
already a scalar term: predecessor is supplied by
Appendix~\ref{sec:hybrid-bootstrap}, the selector is
$C(n,a,b)=I(n,a;u.b)$, addition is $x+y=I(y,x;u.Su)$, and
truncated subtraction is $d(x,y)=I(y,x;u.Pu)$. All other functions below
are finite compositions of these. Thus this proof does not presuppose
that arbitrary primitive recursion has already been compiled. The QFI
conclusion first applies to the scalar language; the general PR compiler
comes afterward.  Literal-equation
induction, together with the primitive-recursive defining equations,
predecessor coverage, and separation of $0$ and $1$, proves induction for every
quantifier-free formula.  In particular, the proof does not assume a faithful
numerical equality test, stability or decidability of equality, successor
uniqueness, or any commutation law for an iterator.  We construct the equality
test and prove its reflection property for arbitrary elements of the
numerical domain.

\subsection{The exact hypotheses}

Work in ordinary intuitionistic first-order logic with equality and genuine
falsity $\bot$.  The language has one numerical sort and a symbol for each
fixed, finite primitive-recursive description.  Distinct descriptions are
not identified merely because they compute the same standard function.
Let $\mathsf{PRE}$ contain the defining equations for zero, successor,
projections, composition, and primitive recursion.  In particular, choose
descriptions $P$ and $C$ with equations
\begin{align}
 P(0)&=0, & P(Sx)&=x, \label{eq:eic-predecessor}\\
 C(0,a,b)&=a, & C(Sx,a,b)&=b. \label{eq:eic-selector}
\end{align}
Write $1=S0$.  In addition to $\mathsf{PRE}$, assume
\begin{align}
 \mathsf{Cov}:&\qquad
       \forall x\,\bigl(x=0\lor x=S(Px)\bigr),
       \label{eq:eic-coverage}\\
 \mathsf{Sep}:&\qquad 0=1\longrightarrow\bot,
       \label{eq:eic-separation}
\end{align}
and induction for each \emph{literal equation} $E(n,\vec p)$:
\begin{equation}
 \forall\vec p\,
 \left[
   \left(E(0,\vec p)\land
       \forall n\bigl(E(n,\vec p)\longrightarrow E(Sn,\vec p)\bigr)
   \right)
   \longrightarrow\forall n\,E(n,\vec p)
 \right].
 \label{eq:eic-induction}
\end{equation}
Here the schematic expression $E$ is a single equality of numerical
terms; the displayed universal quantifiers range only over numbers.
Every description and every instance of the schema is finite external
syntax.

The axiom formulation \eqref{eq:eic-induction} and the usual contextual
induction rule are interchangeable.  For example, if
\[
 \Gamma\vdash E(0,\vec p),
 \qquad
 \Gamma,E(z,\vec p)\vdash E(Sz,\vec p),
\]
where $z$ is absent from $\Gamma$ and the passive parameters, discharge the
second premise's equation and universally introduce $z$.  The axiom then
gives $\Gamma\vdash\forall n E(n,\vec p)$, which can be instantiated at any
target term.  Conversely, assume the antecedent of the axiom and apply the
contextual rule with fresh variables.  This uses the ordinary intuitionistic
logical rules and permits arbitrary, including quantified, formulas in
$\Gamma$.

\subsection{Constructor separation and two elementary tests}

Define the sign description and the countdown description by
\begin{equation}
 \operatorname{sg}(u)=C(u,0,1),
 \qquad
 K(u)=C(u,1,Pu).
 \label{eq:eic-sign-countdown}
\end{equation}
Their defining equations imply
\begin{align}
 \operatorname{sg}(0)&=0,
 &\operatorname{sg}(Sx)&=1,\label{eq:eic-sign-computation}\\
 K(0)&=1,
 &K(Sx)&=x.\label{eq:eic-countdown-computation}
\end{align}

\paragraph{Constructor laws.}
Successor is injective: applying $P$ to $Sx=Sy$ gives $x=y$ by
\eqref{eq:eic-predecessor}.  Successor never equals zero.  Indeed,
$Sx=0$ gives $1=0$ by applying $\operatorname{sg}$ and using
\eqref{eq:eic-sign-computation}; symmetry and
\eqref{eq:eic-separation} give $\bot$.  Consequently $1\ne0$, where
$A\ne B$ abbreviates $A=B\to\bot$.  None of these arguments uses
induction.

\paragraph{The sign test.}
For every $u$,
\begin{align}
 &\operatorname{sg}(u)=0\lor\operatorname{sg}(u)=1,
       \label{eq:eic-sign-range}\\
 &\operatorname{sg}(u)=0\quad\longleftrightarrow\quad u=0,
       \label{eq:eic-sign-zero}\\
 &\operatorname{sg}(u)=1\quad\longleftrightarrow\quad u\ne0.
       \label{eq:eic-sign-one}
\end{align}
To prove the range assertion, split \eqref{eq:eic-coverage}.  The branch
$u=0$ gives sign $0$, and the branch $u=S(Pu)$ gives sign $1$.
The implication $u=0\to\operatorname{sg}(u)=0$ is computation.
For its converse, the successor branch of coverage would make
$\operatorname{sg}(u)=0$ imply $1=0$, a contradiction; in that branch
explosion gives $u=0$ as well.  Similarly, if
$\operatorname{sg}(u)=1$ and $u=0$, computation gives $0=1$.
Conversely, $u\ne0$ excludes the zero branch of coverage, and the
successor branch computes the sign as $1$.

Coverage and \eqref{eq:eic-selector} also give, for all $u,v$,
\begin{equation}
 C(u,v,v)=v.
 \label{eq:eic-selector-idempotence}
\end{equation}
Both coverage branches compute to $v$, so this identity again needs no
induction.

\paragraph{The countdown test and the notation $W$.}
For every $u$,
\begin{equation}
 K(u)=0\quad\longleftrightarrow\quad u=1.
 \label{eq:eic-countdown-zero}
\end{equation}
If $u=1=S0$, \eqref{eq:eic-countdown-computation} gives $K(u)=0$.
Conversely, assume $K(u)=0$ and split coverage.  If $u=0$, then
$K(u)=1$, contradicting separation, and explosion gives $u=1$.
If $u=S(Pu)$, computation gives $K(u)=Pu$; hence $Pu=0$ and
$u=S0=1$.

If a presentation uses a description named $W$ with defining equations
\[
 W(0)=1,\qquad W(Sx)=x,
\]
then it proves $K(u)=W(u)$ for every $u$ already from coverage.
In the zero branch both sides are $1$; in the branch $u=S(Pu)$ both
sides are $Pu$.  Thus $K$ and $W$ are provably equal on the whole domain,
even if they are distinct descriptions.  This equality does not appeal to
standard-model agreement or to uniqueness of an arbitrary recurrence.

\subsection{Four elementary uses of equation induction}

Choose descriptions for addition and truncated subtraction, with recursion
on the second argument:
\begin{align}
 x+0&=x, &x+Sy&=S(x+y),\label{eq:eic-addition}\\
 d(x,0)&=x, &d(x,Sy)&=P(d(x,y)).\label{eq:eic-subtraction}
\end{align}
Using these equations and \eqref{eq:eic-induction}, we now derive
\begin{align}
 0+y&=y,\label{eq:eic-left-zero}\\
 (Sx)+y&=S(x+y),\label{eq:eic-left-successor}\\
 d(Sx,Sy)&=d(x,y),\label{eq:eic-cancel-successors}\\
 d(Sa,a)&=1.\label{eq:eic-diagonal}
\end{align}
Each induction below has a single equation as its induction formula.

For \eqref{eq:eic-left-zero}, induct on $y$.  The base is $0+0=0$.
The step is $0+Sy=S(0+y)=Sy$.

For \eqref{eq:eic-left-successor}, hold $x$ fixed and induct on $y$.
At zero both sides are $Sx$.  At a successor the induction hypothesis
gives
\[
 (Sx)+Sy=S\bigl((Sx)+y\bigr)
       =S\bigl(S(x+y)\bigr)=S(x+Sy).
\]
Combining \eqref{eq:eic-left-successor} with the defining equation for
addition gives the useful rearrangement
\begin{equation}
 (Sy)+r=y+Sr.
 \label{eq:eic-shift-successor}
\end{equation}
Indeed, both sides are $S(y+r)$.

For \eqref{eq:eic-cancel-successors}, hold $x$ fixed and induct on $y$.
The base computes as
\[
 d(Sx,S0)=P(d(Sx,0))=P(Sx)=x=d(x,0).
\]
Applying $P$ to the induction equation gives the step:
\[
 d(Sx,S(Sy))=P(d(Sx,Sy))=P(d(x,y))=d(x,Sy).
\]

Finally, induct on $a$ for \eqref{eq:eic-diagonal}.  The base is
$d(S0,0)=S0=1$.  At a successor,
\[
 d(S(Sa),Sa)=d(Sa,a)=1,
\]
using \eqref{eq:eic-cancel-successors} and the induction hypothesis.

\subsection{An equation invariant that reflects actual equality}

The remaining arithmetic ingredient is the following equation, valid for
all $x,y$:
\begin{equation}
 C\bigl(\operatorname{sg}(d(x,y)),\ x,\ y+d(x,y)\bigr)=x.
 \label{eq:eic-residual-invariant}
\end{equation}
Its meaning is that a nonzero subtraction residual reconstructs the
original argument by addition.  The selector expresses this assertion as
one literal equation, so equation induction applies directly.

Hold $x$ fixed and induct on $y$.  At $y=0$, subtraction gives residual
$x$, and \eqref{eq:eic-left-zero} gives $0+x=x$.  The desired equation
is consequently $C(\operatorname{sg}(x),x,x)=x$, which is
\eqref{eq:eic-selector-idempotence}.

For the step, assume \eqref{eq:eic-residual-invariant} at $y$ and
abbreviate
\[
 r=d(x,y),\qquad r'=Pr=d(x,Sy).
\]
These are local term abbreviations, not additional function symbols.
Split coverage for $r'$.

If $r'=0$, then $\operatorname{sg}(r')=0$, and the successor instance
of \eqref{eq:eic-residual-invariant} reduces immediately to $x=x$.

In the other branch, $r'=S(Pr')$, so constructor separation gives
$r'\ne0$.  Split coverage for $r$.  Its zero branch would imply
$r'=P0=0$, a contradiction, from which the desired step equation follows.
In its successor branch,
\[
 r=S(Pr)=Sr',\qquad \operatorname{sg}(r)=1.
\]
The induction equation therefore computes to $y+r=x$.  Meanwhile,
\eqref{eq:eic-shift-successor} gives
\[
 (Sy)+r'=y+Sr'=y+r=x.
\]
Since $r'$ is a successor, $\operatorname{sg}(r')=1$, so this is
exactly the branch selected in the successor instance of
\eqref{eq:eic-residual-invariant}.  This completes the equation induction.
All case distinctions came from coverage; no induction on an implication
or a disjunction was used.

\paragraph{Reflection at residual $1$.}
We have, for arbitrary $a,b$,
\begin{equation}
 d(Sa,b)=1\quad\longleftrightarrow\quad a=b.
 \label{eq:eic-residual-reflection}
\end{equation}
For the forward direction, substitute $x=Sa$, $y=b$ into
\eqref{eq:eic-residual-invariant}.  If $d(Sa,b)=1$, the sign is $1$
and the selector returns its second branch, giving
\[
 b+1=Sa.
\]
The defining equation for addition gives $b+1=S(b+0)=Sb$.  Thus
$Sb=Sa$, and applying $P$ gives $b=a$, hence $a=b$.
Conversely, $a=b$ and \eqref{eq:eic-diagonal} give
$d(Sa,b)=d(Sa,a)=1$.

Define the equality characteristic by the fixed description
\begin{equation}
 e(a,b)=\operatorname{sg}\bigl(K(d(Sa,b))\bigr).
 \label{eq:eic-equality-characteristic}
\end{equation}
The sign, countdown, and residual reflection lemmas give
\begin{align}
 e(a,b)=0&\quad\longleftrightarrow\quad a=b,
       \label{eq:eic-equality-reflection}\\
 e(a,b)=0&\ \lor\ e(a,b)=1.
       \label{eq:eic-equality-range}
\end{align}
In detail, $e(a,b)=0$ is equivalent to $K(d(Sa,b))=0$ by
\eqref{eq:eic-sign-zero}, which is equivalent to $d(Sa,b)=1$ by
\eqref{eq:eic-countdown-zero}, and then to $a=b$ by
\eqref{eq:eic-residual-reflection}.  The range assertion is
\eqref{eq:eic-sign-range} applied to $K(d(Sa,b))$.

The reflection property concerns the logic's actual equality throughout
the domain.  It was not assumed as part of a numerical comparison symbol,
and its proof did not restrict $a$ or $b$ to standard numerals.
It also yields decidable equality.  In the $e(a,b)=0$ branch,
\eqref{eq:eic-equality-reflection} gives $a=b$.  In the $e(a,b)=1$
branch, an assumption $a=b$ would give $e(a,b)=0$ and therefore $0=1$,
contradicting separation.  Hence
\begin{equation}
 \forall a\,\forall b\,\bigl(a=b\lor(a=b\to\bot)\bigr)
 \label{eq:eic-equality-decidable}
\end{equation}
is a theorem, not an extra hypothesis.

\subsection{Compilation of all quantifier-free formulas}

For each fixed quantifier-free formula $A(\vec p)$, construct a term
$b_A(\vec p)$ by external recursion on its finite syntax.  Take equality,
$\bot$, $\land$, $\lor$, and $\to$ as the primitive formula constructors;
negation is implication into $\bot$.  Set
\begin{align}
 b_{s=t}&=e(s,t), & b_{\bot}&=1,\label{eq:eic-boolean-atoms}\\
 b_{A\land D}&=C(b_A,b_D,1),\label{eq:eic-boolean-and}\\
 b_{A\lor D}&=C(b_A,0,b_D),\label{eq:eic-boolean-or}\\
 b_{A\to D}&=C(b_A,b_D,0).\label{eq:eic-boolean-implies}
\end{align}
Every term here is a finite composition of fixed primitive-recursive
descriptions; if desired, projections and composition give it a single
description symbol.  No internal quantification over formulas or programs
is involved.

An external induction on the syntax of $A$ proves, within the arithmetic
theory,
\begin{equation}
 \bigl(b_A=0\lor b_A=1\bigr)
 \quad\text{and}\quad
 \bigl(b_A=0\longleftrightarrow A\bigr).
 \label{eq:eic-boolean-correctness}
\end{equation}
Here are the full connective checks.  In every case $b_A=1$ implies
$\neg A$: if $A$ held, the already established implication
$A\to b_A=0$ would contradict separation.

\paragraph{Atomic formulas and falsity.}
For an equation, both claims are
\eqref{eq:eic-equality-reflection}--\eqref{eq:eic-equality-range}.
For falsity the range assertion is $1=0\lor1=1$, with its right disjunct
proved by reflexivity.  Also $1=0\to\bot$ follows from separation and
symmetry, while $\bot\to1=0$ follows by explosion.  Thus $1=0$ is
equivalent to $\bot$, as required.

\paragraph{Conjunction.}
Split $b_A=0\lor b_A=1$.  In the zero branch, $A$ holds and
$b_{A\land D}=b_D$.  Its range follows from the range of $b_D$,
and its being zero is equivalent to $D$, hence to $A\land D$ under
the known $A$.  In the one branch, $\neg A$ holds and
$b_{A\land D}=1$.  Its range is immediate.  If it were zero,
separation would give $\bot$, hence $A\land D$; conversely an
assumption $A\land D$ contradicts $\neg A$, and explosion gives
$b_{A\land D}=0$.

\paragraph{Disjunction.}
In the branch $b_A=0$, the formula $A$ holds and
$b_{A\lor D}=0$.  The range is immediate, and both $b_{A\lor D}=0$
and $A\lor D$ hold, proving their equivalence.  In the branch
$b_A=1$, we have $\neg A$ and $b_{A\lor D}=b_D$.  Its range follows
from the range of $b_D$.  If its value is zero, then $D$ and hence
$A\lor D$.  Conversely, from $A\lor D$, the $A$ alternative
contradicts $\neg A$ and yields $b_D=0$ by explosion, while the
$D$ alternative yields $b_D=0$ directly.  This proves the required
equivalence in both alternatives.

\paragraph{Implication.}
In the branch $b_A=0$, the formula $A$ holds and
$b_{A\to D}=b_D$.  Its range follows from the range of $b_D$.
If its value is zero, then $D$, which gives $A\to D$.  Conversely,
$A\to D$ and the known $A$ give $D$, hence $b_D=0$.
In the branch $b_A=1$, we have $\neg A$ and
$b_{A\to D}=0$.  Its range is immediate.  To prove $A\to D$,
assume $A$; this contradicts $\neg A$, and explosion gives $D$.
Thus both $b_{A\to D}=0$ and $A\to D$ hold in this branch, and
their equivalence follows.  This argument uses the proved two-value
range of the finite characteristic.  It does not use excluded middle
for arbitrary formulas or change the surrounding intuitionistic logic.

This completes the external syntactic induction proving
\eqref{eq:eic-boolean-correctness} for every quantifier-free formula.

\subsection{Recovery of induction and the full theory}

Let $A(n,\vec p)$ be any fixed quantifier-free formula.  Suppose an
arbitrary retained context $\Gamma$ proves its base and its induction
step:
\[
 \Gamma\vdash A(0,\vec p),
 \qquad
 \Gamma,A(z,\vec p)\vdash A(Sz,\vec p),
\]
with the usual freshness condition on $z$.  By
\eqref{eq:eic-boolean-correctness}, the base gives
$\Gamma\vdash b_A(0,\vec p)=0$.  Under the additional assumption
$b_A(z,\vec p)=0$, the same equivalence gives $A(z,\vec p)$,
the given step gives $A(Sz,\vec p)$, and the equivalence gives
$b_A(Sz,\vec p)=0$.  Apply literal-equation induction to
$b_A(n,\vec p)=0$ and then apply its reflection implication.
The result is $\Gamma\vdash A(t,\vec p)$ for every target term $t$.
Thus the theory proves the full contextual quantifier-free induction
rule, including in arbitrary quantified contexts.

\paragraph{Completion theorem.}
In the common language of fixed primitive-recursive descriptions, the
theory
\[
 \mathsf{PRE}+\mathsf{Cov}+\mathsf{Sep}+\mathsf{EI}
\]
has exactly the same sequents as ordinary first-order intuitionistic
primitive recursive arithmetic with all quantifier-free induction
instances.  The forward simulation follows from the recovered
quantifier-free induction, successor injectivity, and successor separation;
all primitive-recursive defining equations and all logical rules are
shared.  For the reverse simulation, literal-equation induction is an
instance of quantifier-free induction, and successor separation gives
$\mathsf{Sep}$.  Quantifier-free induction proves $\mathsf{Cov}$:
at zero choose its left disjunct, and at a successor use
$Sx=S(P(Sx))$, which follows from $P(Sx)=x$, for its right disjunct.

Both interpretations therefore use the full domain, defined by $x=x$,
and preserve every function symbol, equality, logical connective, and
quantifier.  Both composites have the identity isomorphism, with graph
$x=y$.  The theories are consequently identity-domain bi-interpretable
and equiconsistent.  This completion theorem is unconditional once the
literal-equation induction rule has been obtained; it has no dependency
on a classical successor-uniqueness reduction.

\clearpage
\section{Pairing, primitive recursion, and internal numerical composites}
\label{sec:iterator-coding}
This appendix is adapted with attribution from the progress-closure
arithmetic development~\cite{progress}. Its prerequisites have now been
proved: scalar predecessor, arithmetic reflection, and full quantifier-free
induction. Pairing is a derived proof tool; it is not part of any checker.

\subsection{Pairing from a bounded square-root search}

Once QFI is available, a familiar square-based injection gives a
particularly direct compiler. Write $a\le b$ for $d(a,b)=0$.
The usual linear-order and semiring identities follow by fixed
quantifier-free inductions from the successor, addition, and monus
equations. Multiplication itself is the scalar term
\[
                        a\cdot b=I(b,0;u.u+a).
\]
Only primitive-recursive arithmetic is involved in the following
bounded root search; there is no minimization over an unbounded domain.

For a fixed parameter $n$, put
\[
 U_n(v)=C(d((Sv)^2,n),Sv,v),\qquad
 R_n(i)=I(i,0;v.U_n(v)),\qquad q(n)=R_n(n).
\]
The update increments $v$ exactly when $(v+1)^2\le n$, and
otherwise retains it. QFI in $i$, with $n$ fixed, proves
\[
 R_n(i)\le i,\qquad R_n(i)^2\le n,\qquad
 R_n(i)=i\ \vee\ n<(S(R_n(i)))^2.              \tag{Root invariant}
\]
The base has $R_n(0)=0$. If the update increments, its test
contradicts the strict-upper-bound alternative in the third invariant.
Thus the old state equals $i$, and the new state equals $Si$;
its square bound is the increment test. If the update does not increment, the
strict upper square bound holds and the state remains fixed.
An already established strict upper square bound forces every
subsequent update to remain fixed. These are exactly the cases
needed for the three displayed invariants.

At $i=n$, either the strict upper square bound already holds or
$R_n(n)=n$, in which case $n<(n+1)^2$ supplies it. Thus
\[
                        q(n)^2\le n<(q(n)+1)^2.           \tag{Root}
\]
Monotonicity of squares shows uniqueness of a number satisfying
Root. All formulas just used are quantifier free in the scalar
language, with fixed $n$ during the induction.

Now define
\[
 \pi(a,b)=(a+b)^2+a,\qquad
 L(n)=d(n,q(n)^2),\qquad H(n)=d(q(n),L(n)).
\]
For $k=a+b$, one has $k^2\le k^2+a<(k+1)^2$ since $a\le k$.
Root uniqueness gives $q(\pi(a,b))=a+b$, and therefore
\[
                    L(\pi(a,b))=a,\qquad H(\pi(a,b))=b.
\]
Surjectivity of $\pi$ is unnecessary. The original numerical domain
is never replaced by the range of this coding; pairing is used only
to store the finite state of a fixed primitive recursion. This
standard square pairing is a convenient reconstruction here, not
a claimed new arithmetic basis.

\subsection{Compiling each primitive recursion}

Suppose the fixed target description is
\[
 f(0,\vec p)=g(\vec p),\qquad
 f(Sn,\vec p)=h(n,f(n,\vec p),\vec p).
\]
Assume by external structural induction that the earlier $g,h$
already have scalar translations $g^*,h^*$. Define
\[
 W(u,\vec p)=\pi(SL(u),h^*(L(u),H(u),\vec p)),
\]
\[
 V(n,\vec p)=I(n,\pi(0,g^*(\vec p));u.W(u,\vec p)),
 \qquad f^*(n,\vec p)=H(V(n,\vec p)).
\]
The update has only its state variable bound, with every $p_i$
fixed. Equation induction proves $L(V(n,\vec p))=n$.
The two projection equations and scalar computation then prove
both defining equations for $f^*$ internally. Initial functions
and composition translate directly. This completes the translation
of every fixed finite PR description without adding a universal
evaluator.

\subsection{Why both numerical round trips are internal}

Two fixed scalar iterators with the same seed and provably equal
updates are internally equal at all counters: use literal-equation
induction comparing their values. Their base agrees, and their
successor equations plus congruence prove the step. The same
uniqueness argument works for ordinary primitive recursion.

For an ordinary PR description, induct externally on its finite
construction. The double translation satisfies its original
defining equations, so internal equation induction identifies it
with the original operation. For a native scalar description,
earlier translated updates agree by the external induction
hypothesis; scalar uniqueness identifies the double-translated
iterator. These proofs hold at arbitrary internal counters.

Thus the domain relation $x=y$ is provably a total bijection
preserving every fixed numerical operation in both composites.
The formula constructions and round trips in
Section~\ref{sec:choice-introduction} give the corresponding formula
composites. Both interpretations preserve falsity, taking native
$0=1$ to an ordinary contradiction and ordinary $\bot$ to native
$0=1$. Finite proof translations therefore establish equiconsistency.
\clearpage
\section{Complete primary verifier}

The following listing is the complete primary acceptance module.
The normalization and source-size procedures do not remove executable
code from it.

\begin{lstlisting}[language=Python]
"""Five-rule clausal iPRA with zero-normal iterators and uniform successor Step.

Numerical syntax: 0; integer bound indices (>0) or free levels (<0);
['s', term]; ['i', counter, seed, body]. Only body binds one variable.
Formula syntax: ['=', left, right]; ['d', matrix]. D binds x then y,
so the matrix uses index 1 for y and 2 for x. Zero-counter iterators
are excluded, and opening contracts them when substitution exposes one.
"""
import json
import sys


def require(condition):
    if not condition:
        raise ValueError('Invalid certificate')


def iterator(n, a, h):
    return a if n == 0 else ['i', n, a, h]


def term(a, scope, depth=0):
    if type(a) is int:
        return -scope <= a <= depth
    if type(a) is not list:
        return False
    if len(a) == 2 and a[0] == 's':
        return term(a[1], scope, depth)
    return (len(a) == 4 and a[0] == 'i' and a[1] != 0
            and term(a[1], scope, depth) and term(a[2], scope, depth)
            and term(a[3], scope, depth + 1))


def formula(a, scope, depth=0):
    return (type(a) is list
            and (len(a) == 3 and a[0] == '='
                 and all(term(t, scope, depth) for t in a[1:])
                 or len(a) == 2 and a[0] == 'd'
                 and formula(a[1], scope, depth + 2)))


def opening(a, value, depth=0):
    if type(a) is int:
        return value if a == depth + 1 else a - (a > depth + 1)
    tag = a[0]
    b = [opening(t, value, depth + (2 if tag == 'd' else int(tag == 'i' and i == 2)))
         for i, t in enumerate(a[1:])]
    return iterator(*b) if tag == 'i' else [tag] + b


def derive(proof, context, scope):
    match proof:
        case ['u', index, witness, premise]:
            require(type(index) is int and 0 <= index < len(context)
                    and term(witness, scope))
            antecedent, body = context[index]
            require(derive(premise, context, scope) == antecedent)
            return opening(body, witness)
        case ['r', a]:
            require(term(a, scope))
            return ['=', a, a]
        case ['x', major, minor]:
            match derive(major, context, scope), derive(minor, context, scope):
                case ['=', s, t], ['=', ['i', ['s', n], a, h], r]:
                    require(s == iterator(n, a, h))
                    return ['=', r, opening(h, t)]
        case [('e' | 'i') as tag, goal, data, branch]:
            require(formula(goal, scope))
            fresh = -scope - 1
            if tag == 'e':
                paired = derive(data, context, scope)
                require(paired[0] == 'd')
                body = opening(paired[1], fresh, 1)
                clause = (opening(body, 0), body)
                target = goal
            else:
                require(goal[0] == 'd' and term(data, scope))
                body = opening(goal[1], data, 1)
                eta = fresh if goal[1][0] == '=' and data == 0 else 0
                clause = (['=', 0, 0], opening(body, eta))
                target = opening(body, ['s', fresh])
            require(derive(branch, context + [clause], scope + 1) == target)
            return goal
    raise ValueError('Invalid proof')


def verify(data):
    require(type(data) is dict and set(data) == {'goal', 'proof'})
    require(formula(data['goal'], 0))
    require(derive(data['proof'], [], 0) == data['goal'])
    return True


if __name__ == '__main__':
    try:
        verify(json.load(sys.stdin))
    except Exception:
        sys.exit(1)
    print('ACCEPT')
\end{lstlisting}

\clearpage
\begin{thebibliography}{9}
\raggedright
\bibitem{baseline}
Vladimir Reshetnikov, \emph{ProveIt: Concise presentations of
intuitionistic primitive recursive arithmetic}, commit
\texttt{6ca328039cafa1537d81ee6b67ee2dc54ffa5a4e}.
\href{https://github.com/VladimirReshetnikov/ProveIt/tree/6ca328039cafa1537d81ee6b67ee2dc54ffa5a4e/Logic/PrimitiveRecursiveArithmetic/ConciseIPRA}{ConciseIPRA at the frozen comparison commit}.

\bibitem{erasure}
ProveIt research package, \emph{Internal scalar predecessor bootstrap},
\path{Research/certified_zero/reference/scalar_predecessor.tex},
at the commit in~\cite{baseline}. Reproduced with attribution.

\bibitem{forward}
ProveIt research package, \emph{From literal-equation induction to
full intuitionistic PRA},
\path{Research/certified_zero/reference/equality_reflection.tex},
at the commit in~\cite{baseline}. Reproduced with attribution.

\bibitem{progress}
ProveIt research package, \emph{Pairing, primitive recursion, and
internal numerical composites},
\path{Research/certified_zero/reference/primitive_recursion.tex},
at the commit in~\cite{baseline}. Reproduced with attribution.

\bibitem{choice}
ProveIt research package, \emph{Complementary introductions with
one local assumption},
\path{Research/certified_zero/reference/inherited_choice_basis.tex},
and \emph{Arithmetic strength and both ordinary interpretation
composites},
\texttt{inherited\_ordinary\_bridge.tex},
at the commit in~\cite{baseline}. Reproduced with attribution.

\bibitem{zero}
ProveIt, \emph{ZeroNormalization} and \emph{NativeZeroNormal},
Lean sources under \path{Lean/ConciseIPRA/}, at the commit
in~\cite{baseline}.

\bibitem{eval}
ProveIt, \emph{Computational equality: six rules},
\path{Research/computational_equality/} and
\path{Lean/ConciseIPRA/NativeEval.lean}, at the commit
in~\cite{baseline}.

\bibitem{incoming}
ProveIt incoming archives,
\path{docs/incoming/ipra_clausal_canonical.zip} and
\path{docs/incoming/ipra_scalar_uniqueness.zip},
at the commit in~\cite{baseline}. The former supplies the incoming
five-rule comparison; the latter supplies the 866-byte profile
mentioned only for scope of comparison.
\end{thebibliography}
\end{document}
